\documentclass[11pt,letterpaper]{article}
\usepackage[utf8]{inputenc}
\usepackage[T1]{fontenc}
\usepackage[spanish,es-nodecimaldot]{babel}
\usepackage{lmodern}
\usepackage{microtype}
\usepackage{geometry}
\geometry{left=2.5cm,right=2.2cm,top=2.4cm,bottom=2.4cm}
\usepackage{amsmath,amssymb,mathtools,bm}
\usepackage{siunitx}
\sisetup{output-decimal-marker={,},per-mode=symbol}
\usepackage{booktabs,longtable,tabularx,array}
\usepackage{xcolor}
\definecolor{novaBlue}{HTML}{1F4E79}
\definecolor{novaLight}{HTML}{EAF2F8}
\usepackage{enumitem}
\usepackage{tcolorbox}
\tcbuselibrary{breakable}
\usepackage{xurl}
\usepackage{hyperref}
\hypersetup{colorlinks=true,linkcolor=novaBlue,urlcolor=novaBlue,citecolor=novaBlue,
 pdftitle={Modelo matemático Nova Spot Micro}}
\usepackage{cleveref}
\usepackage{fancyhdr}
\pagestyle{fancy}\fancyhf{}
\fancyhead[L]{Modelo matemático Nova Spot Micro}
\fancyhead[R]{NovaSM3}\fancyfoot[C]{\thepage}
\setlength{\headheight}{14pt}
\setcounter{tocdepth}{2}
\newcommand{\vect}[1]{\bm{#1}}
\newcommand{\mat}[1]{\bm{#1}}
\newcommand{\R}{\mathbb{R}}
\newcommand{\trans}{^{\mathsf T}}
\newcommand{\sat}{\operatorname{sat}}
\newcommand{\clip}{\operatorname{clip}}
\newcommand{\diag}{\operatorname{diag}}
\newcommand{\rank}{\operatorname{rank}}
\newcommand{\sign}{\operatorname{sign}}
\newcommand{\FK}{\operatorname{FK}}
\newcommand{\IK}{\operatorname{IK}}
\newcommand{\argmin}{\operatorname*{arg\,min}}
\newcolumntype{Y}{>{\raggedright\arraybackslash}X}
\newtcolorbox{validity}{colback=novaLight,colframe=novaBlue,
 title=Estado de validez,fonttitle=\bfseries}
\newtcolorbox{workedexample}[1]{breakable,colback=white,colframe=novaBlue,
 title={Ejemplo desarrollado: #1},fonttitle=\bfseries,
 before skip=8pt,after skip=8pt}
\newtcolorbox{projectresult}[1]{breakable,colback=novaLight,colframe=novaBlue,
 title={Resultado para Nova: #1},fonttitle=\bfseries,
 before skip=8pt,after skip=8pt}
\newtcolorbox{tracebox}[1]{breakable,colback=gray!4,colframe=gray!65,
 title={Trazabilidad: #1},fonttitle=\bfseries,boxrule=0.6pt,
 before skip=6pt,after skip=6pt}

\begin{document}
\hypersetup{pageanchor=false}
\pagenumbering{roman}
\begin{titlepage}
\centering
\vspace*{3cm}
{\Huge\bfseries\color{novaBlue} Modelo matemático del robot cuadrúpedo\\[0.3cm]
Nova Spot Micro\par}
\vspace{2cm}
{\Large Modelo nominal computable para postura, paso, gateo y otros patrones de locomoción\par}
\vfill
{\large Proyecto de tesis --- NovaSM3\par}
\vspace{1cm}
ROS 2 Humble $\cdot$ Gazebo Sim $\cdot$ MuJoCo
\vfill
Versión LaTeX ampliada con control discreto: 16 de agosto de 2026
\end{titlepage}

\tableofcontents
\clearpage
\hypersetup{pageanchor=true}
\pagenumbering{arabic}

\section*{Resumen}\addcontentsline{toc}{section}{Resumen}
Este documento presenta el modelo matemático nominal del robot cuadrúpedo Nova
Spot Micro. Integra geometría, cinemática directa e inversa, espacio de trabajo,
Jacobianos, singularidades, cinemática de segundo orden, dinámica rígida,
actuadores, contacto pie--suelo, centro de masa, estabilidad, generación de
marchas, teoría y diseño de control discreto, identificación y validación. Su estructura es independiente del patrón
de locomoción y puede aplicarse a postura, paso, gateo y otros movimientos que
respeten las hipótesis y límites declarados.

\begin{validity}
La estructura matemática es computable y ha sido verificada internamente. Las
dimensiones generales y la masa total tienen respaldo público; la distribución
de masa, inercias, fricción, contacto y actuadores contiene valores nominales de
ingeniería hasta identificar el ejemplar físico. En consecuencia, este es un
\emph{modelo nominal completo y computable}, no un gemelo digital validado.
\end{validity}

\section{Alcance, hipótesis y niveles de modelado}
El modelo completo emplea una base flotante de seis grados de libertad, doce
coordenadas articulares y fuerzas de contacto. El modelo reducido representa
una pata como cadena serial de tres grados de libertad. El primero sirve para
pose corporal, contactos y estabilidad global; el segundo para FK, IK,
Jacobianos, workspace, pares locales y trayectorias cartesianas.
\begin{align}
 \vect q &= [\vect r_B\trans,\,\vect\eta_B\trans,\,\vect q_J\trans]\trans
 \in\R^{18}, &
 \vect q_i &= [q_{c,i},q_{f,i},q_{t,i}]\trans\in\R^3. \label{eq:qlevels}
\end{align}
Se asumen eslabones rígidos, suelo nominal plano, contacto unilateral, control
de posición en los servos y gateo suficientemente lento para aplicar el
criterio cuasiestático. Holgura, flexibilidad y dispersión de fabricación se
representarán mediante incertidumbre. Una coincidencia entre ecuaciones y
simuladores constituye verificación de implementación, no validación física.

\section{Convenciones, marcos y nomenclatura}
Se usa el Sistema Internacional y REP--103: $x$ hacia delante, $y$ hacia la
izquierda y $z$ hacia arriba. Los marcos son $\{W\}$ (mundo), $\{B\}$ (cuerpo),
$\{H_i\}$ (cadera), $\{C_i\}$ (coxa), $\{F_i\}$ (fémur), $\{T_i\}$ (tibia) y
$\{P_i\}$ (pie). FL, FR, RL y RR designan las cuatro patas. El factor
$s_i=+1$ para patas izquierdas y $s_i=-1$ para derechas.
\begin{table}[ht]
\centering\caption{Posición de los pivotes de coxa en el marco corporal.}
\begin{tabular}{ccc}\toprule
Pata & $s_i$ & ${}^B\vect p_{H_i}$ \\ \midrule
FL & $+1$ & $[d_x/2,d_y/2,0]\trans$ \\
FR & $-1$ & $[d_x/2,-d_y/2,0]\trans$ \\
RL & $+1$ & $[-d_x/2,d_y/2,0]\trans$ \\
RR & $-1$ & $[-d_x/2,-d_y/2,0]\trans$ \\ \bottomrule
\end{tabular}\end{table}

\section{Cómo se construyó el modelo y cómo leer su trazabilidad}
El modelado no se obtuvo de una sola fórmula ni de copiar un robot distinto. Se
siguió el método de modelado físico por capas:
\begin{enumerate}
 \item definir marcos, signos y unidades;
 \item obtener geometría y parámetros nominales del NovaSM3;
 \item derivar la cadena cinemática bajo esas convenciones;
 \item comprobar FK e IK entre sí y contra el URDF;
 \item construir Jacobianos y energía para obtener dinámica;
 \item añadir actuador, contacto y estabilidad con supuestos declarados;
 \item discretizar únicamente después de fijar el modelo continuo y el periodo;
 \item contrastar cada salida calculada con pruebas o ensayos independientes.
\end{enumerate}

Cada familia de ecuaciones tiene cuatro niveles de trazabilidad:
\begin{description}[style=nextline,leftmargin=2.4cm]
 \item[Base T] fuente teórica u oficial que justifica el método general;
 \item[Adaptación A] decisiones propias de ejes, signos, dimensiones y alcance;
 \item[Código C] archivo o función que evalúa la formulación;
 \item[Evidencia E] prueba automatizada, simulación o ensayo que la contrasta.
\end{description}
Una fuente teórica no valida el robot, y un ensayo favorable no demuestra por sí
solo que la ecuación sea correcta. La tesis exige conservar los cuatro niveles.

\subsection{Fuentes base identificadas}
\small
\begin{longtable}{p{0.7cm}p{4.2cm}p{8.0cm}}
\caption{Fuentes de las familias matemáticas.}\\
\toprule ID&Fuente&Uso y límite en esta tesis\\\midrule\endfirsthead
\toprule ID&Fuente&Uso y límite en esta tesis\\\midrule\endhead
T01&REP--103, ROS&Convención $x$ adelante, $y$ izquierda, $z$ arriba. Se usa
para evitar cambios de signo entre modelo, ROS 2 y simuladores.\\
T02&Craig, \emph{Introduction to Robotics}&Transformaciones homogéneas,
cinemática serial y Jacobiano. La cadena concreta y sus signos se derivaron
nuevamente para Nova.\\
T03&Mike4192, \path{spot_micro_kinematics_python}&Referencia conceptual MIT
para geometría de Spot Micro e IK. No se copiaron dimensiones ni parámetros de
marcha; la implementación Nova es propia.\\
T04&Formulación de Euler--Lagrange y dinámica multicuerpo; Featherstone y notas
MIT Underactuated&Forma $M(q)\ddot q+C(q,\dot q)\dot q+g(q)$, positividad de
$M$ y fuerza generalizada $J^T f$.\\
T05&Contacto unilateral de Signorini y fricción de Coulomb&Complementariedad,
cono de fricción y modelo penalizado. El resorte--amortiguador es una
aproximación nominal, no el algoritmo interno exacto de Gazebo o MuJoCo.\\
T06&TowerPro MG996R, documentación del fabricante&Voltaje, velocidad, par de
bloqueo y corrientes de catálogo. No proporciona par continuo identificado.\\
T07&Ogata; Franklin, Powell y Workman&Muestreo, ZOH, transformada Z,
discretización, Jury, lugar de raíces, PID y respuesta en frecuencia.\\
T08&Bellman y formulación LQR; notas MIT Underactuated&Índice cuadrático,
ecuación de Riccati y realimentación óptima para el modelo lineal declarado.\\
T09&Kalman, 1960, DOI 10.1115/1.3662552&Estimación recursiva lineal y evolución
de covarianza. Su inclusión es diseño futuro hasta disponer de IMU y ruido
caracterizado.\\
T10&Documentación oficial \path{trajectory_msgs} y
\path{joint_trajectory_controller}&Contrato de nombres, puntos, posiciones,
velocidades y tiempo usado para enviar referencias desde ROS 2.\\
T11&Documentación oficial de sensores Gazebo&Significado y publicación de
sensores de contacto. Justifica la adquisición, no el plan ideal de apoyo.\\
\bottomrule
\end{longtable}
\normalsize

\subsection{Matriz de trazabilidad de las ecuaciones Nova}
\small
\begin{longtable}{p{0.7cm}p{2.5cm}p{1.7cm}p{3.2cm}p{4.3cm}}
\caption{Cadena teoría--adaptación--código--evidencia.}\\
\toprule ID&Resultado buscado&Base&Implementación&Comprobación y uso en tesis\\
\midrule\endfirsthead
\toprule ID&Resultado buscado&Base&Implementación&Comprobación y uso en tesis\\
\midrule\endhead
A01&Posición 3D del pie&T01--T03&\path{forward_leg}&Ida y vuelta FK--IK;
genera objetivos y permite reportar trayectoria.\\
A02&Ángulos desde posición&T02--T03&\path{inverse_leg}&Dominio, límites,
rama de rodilla y continuidad; convierte marcha cartesiana a 12 referencias.\\
A03&Velocidad y fuerza--par&T02,T04&\path{leg_jacobian},
\path{foot_force_to_joint_torque}&Diferencias centrales y potencia
virtual; cuantifica singularidad y carga articular.\\
A04&Masa, Coriolis y gravedad&T04&\path{leg_mass_matrix},
\path{coriolis_torque}, \path{gravity_torque}&Simetría, autovalores
positivos y casos finitos; sustenta el modelo nominal de esfuerzo.\\
A05&Capacidad del MG996R&T06&\path{actuator_torque_limit},
\path{actuator_current}&Comparación preliminar par requerido/disponible;
dimensionamiento, no autorización de hardware.\\
A06&Contacto y soporte&T05,T11&URDF, \path{contact_monitor},
\path{static_stability_margin}&Sensores Gazebo y bolsa de 76 ciclos;
detectó que el patrón ideal no se cumple.\\
A07&Trayectoria de gateo&T02--T03&\path{cartesian_crawl}&Alcanzabilidad,
continuidad, cadencia y rosbag; línea base nominal de comparación.\\
A08&Marcha paso&T02--T03&\path{cartesian_step_walk}&Dos ensayos Gazebo y
uno MuJoCo; segunda marcha nominal de la tesis.\\
A09&Tiempo discreto&T07,T10&\path{advance_phase_deadline}&Prueba unitaria y
ciclos de \SI{4.32}{s}; evita deriva de fase.\\
A10&Modelo discreto local&T04,T07&Exponencial por bloques en
\path{generar_ejemplos_modelo.py}&Polos y rango de controlabilidad; puente
entre dinámica y futuro diseño digital.\\
A11&Control LQR nominal&T08&\path{solve_discrete_are} en el generador&Polos
cerrados y respuesta a $10^\circ$; ejemplo de diseño, no implementación actual.\\
A12&Retardos de apoyo&T11 y método propio&\path{analizar_contactos_rosbag.py}&
76 transiciones delanteras y ausencia trasera; criterio para corregir gateo.\\
\bottomrule
\end{longtable}
\normalsize

\subsection{Trazabilidad de código y evidencia}
\small
\begin{longtable}{p{0.7cm}p{5.5cm}p{6.7cm}}
\toprule ID&Archivo o evidencia&Qué demuestra exactamente\\\midrule\endfirsthead
\toprule ID&Archivo o evidencia&Qué demuestra exactamente\\\midrule\endhead
C01--C03&\path{src/nova_gait_controller/nova_gait_controller/kinematics.py}&
FK, IK, Jacobiano geométrico indirecto y generadores cartesianos.\\
C04--C06&\path{src/nova_gait_controller/nova_gait_controller/mathematical_model.py}&
Jacobiano analítico, dinámica, actuador, contacto, COM y margen nominal.\\
C07--C09&\path{src/nova_gait_controller/nova_gait_controller/gait_controller.py}&
Máquina de estados, validación, deadline, trayectoria y fase publicada.\\
C10--C11&\path{scripts/generar_ejemplos_modelo.py}&Cálculos numéricos que se
transcriben en las cajas desarrolladas; salida JSON reproducible.\\
C12&\path{Experimentos/analizar_contactos_rosbag.py}&Emparejamiento de
transiciones, retardos y coincidencia ponderada por tiempo.\\
E01--E05&\path{src/nova_gait_controller/test/}&Pruebas de cinemática, modelo,
temporización, seguridad y consistencia entre representaciones.\\
E06&\path{Experimentos/rosbag2/contactos_gateo_validado_20260814_1410}&Bolsa
aceptada de 76 ciclos con fase, contactos, métricas y cero eventos de seguridad.\\
E07&\path{Experimentos/analisis/contactos_gateo_validado_20260814_1410/INFORME_CONTACTOS.md}&
Retardos, porcentajes por pata y 32.550\% de coincidencia simultánea.\\
E08&\path{Experimentos/comparacion_cadencia_corregida_20260814/INFORME_REPRODUCIBILIDAD.md}&
Repetibilidad de la cadencia corregida en dos ensayos independientes.\\
E09&\path{Documentacion/MARCHA_PASO_VALIDACION.md}&Dos ensayos Gazebo y
validación articular de 12 ciclos en MuJoCo.\\
\bottomrule
\end{longtable}
\normalsize

\subsection{Estructura obligatoria de los ejemplos}
Cada ejemplo desarrollado responde, en este orden, siete preguntas:
\begin{enumerate}
 \item \textbf{¿Qué se quiere obtener?} Define la salida útil para la tesis.
 \item \textbf{¿De dónde sale el método?} Identifica T, A, C y E.
 \item \textbf{¿Qué datos entran?} Lista valor, unidad y estado de validez.
 \item \textbf{¿Qué hipótesis se usan?} Expone simplificaciones y dominio.
 \item \textbf{¿Cómo se sustituye?} Conserva pasos y unidades intermedias.
 \item \textbf{¿Cómo se comprueba?} Usa identidad, prueba o evidencia separada.
 \item \textbf{¿Qué significa?} Indica para qué decisión de tesis sirve y qué
 no permite afirmar.
\end{enumerate}
Los ejemplos que usan parámetros nominales son ejercicios de verificación del
modelo. Los que usan rosbag son resultados experimentales de simulación. Ninguno
se etiqueta como resultado físico mientras no exista medición del ejemplar.

\section{Diccionario razonado de variables}
Las variables no se incluyen por costumbre: cada una responde una pregunta de
la tesis. Su estado se clasifica como \emph{entrada}, \emph{parámetro},
\emph{estado}, \emph{salida calculada} o \emph{medición}.

\subsection{Geometría y cinemática}
\footnotesize
\begin{longtable}{p{2.2cm}p{1.7cm}p{3.9cm}p{5.5cm}}
\toprule Símbolo&Unidad o tipo&Qué representa&Por qué se obtiene\\\midrule\endfirsthead
\toprule Símbolo&Unidad o tipo&Qué representa&Por qué se obtiene\\\midrule\endhead
$i$&índice&FL, FR, RL o RR&Permite aplicar una ecuación a cuatro patas sin
duplicarla.\\
$s_i$&$\{-1,+1\}$&Signo izquierda/derecha&Corrige simetría de coxa y evita
ecuaciones manuales con signos incompatibles.\\
$d_x,d_y$&m&Separaciones de caderas&Ubican cada pata respecto al cuerpo y
determinan el polígono de soporte.\\
$l_c,l_f,l_t$&m&Longitudes de coxa, fémur y tibia&Definen alcance, FK, IK,
Jacobiano y dinámica; deben medirse físicamente después.\\
$q_c,q_f,q_t$&rad&Ángulos articulares&Son estados y referencias que puede
transportar ROS 2; conectan modelo cartesiano con actuadores.\\
$\vect p_i=[x,y,z]^T$&m&Posición del pie&Es la variable natural para diseñar
altura, longitud de paso y ubicación de contacto.\\
$Q=q_f+q_t$&rad&Ángulo acumulado planar&Reduce notación y expresa la orientación
de la tibia respecto al fémur.\\
$\mat J_i$&m/rad&Jacobiano del pie&Convierte velocidad articular en cartesiana y
fuerza del pie en par; permite detectar singularidades.\\
$\sigma_{min},\kappa$&m/rad, --&Valor singular mínimo y condición&Cuantifican
pérdida de movilidad y sensibilidad numérica antes de ordenar un objetivo.\\
\bottomrule
\end{longtable}
\normalsize

\subsection{Dinámica, actuador, contacto y estabilidad}
\footnotesize
\begin{longtable}{p{2.2cm}p{1.7cm}p{3.9cm}p{5.5cm}}
\toprule Símbolo&Unidad o tipo&Qué representa&Por qué se obtiene\\\midrule\endfirsthead
\toprule Símbolo&Unidad o tipo&Qué representa&Por qué se obtiene\\\midrule\endhead
$\mat M(q)$&kg m$^2$&Inercia articular&Relaciona aceleración con par y permite
comparar demanda frente al actuador.\\
$\vect c$&N m&Coriolis y centrífugo&Captura acoplamiento dependiente de velocidad;
evita atribuir todo el par a gravedad.\\
$\vect g$&N m&Par gravitacional&Cuantifica el esfuerzo para sostener postura.\\
$\vect\tau_f$&N m&Fricción pasiva&Representa disipación nominal y permite hacer
explícita una incertidumbre importante.\\
$\vect f_i$&N&Fuerza cartesiana del pie&Entra por $J^Tf$ y conecta contacto con
carga articular.\\
$\vect\tau$&N m&Par requerido&Salida central para seguridad, selección de servo
y comparación de controladores.\\
$V,\omega$&V, rad/s&Tensión y velocidad del servo&Modifican el par disponible;
impiden usar el par de bloqueo como constante universal.\\
$I$&A&Corriente estimada o medida&Sirve para dimensionar fuente, energía y
criterios de protección.\\
$\phi,\delta$&m&Separación y penetración&Deciden si existe contacto y cuánto
actúa el modelo penalizado.\\
$\lambda_n,\vect\lambda_t$&N&Fuerza normal y tangencial&Permiten evaluar apoyo,
fricción y deslizamiento.\\
$\vect r_{COM}$&m&Centro de masa total&Se proyecta sobre contactos para medir
estabilidad, no solo inclinación visual.\\
$m_s$&m&Margen estático firmado&Será una métrica de tesis: positivo dentro del
soporte, negativo fuera.\\
\bottomrule
\end{longtable}
\normalsize

\subsection{Marcha, control discreto y experimento}
\footnotesize
\begin{longtable}{p{2.2cm}p{1.7cm}p{3.9cm}p{5.5cm}}
\toprule Símbolo&Unidad o tipo&Qué representa&Por qué se obtiene\\\midrule\endfirsthead
\toprule Símbolo&Unidad o tipo&Qué representa&Por qué se obtiene\\\midrule\endhead
$L,H$&m&Longitud y elevación del paso&Parámetros independientes que determinan
avance, despeje y exigencia articular.\\
$N,\Delta t,T$&--, s, s&Muestras, periodo de referencia y ciclo&Definen la
cadencia reproducible y la segmentación de rosbag.\\
$\beta,\varphi_i$&--&Factor de apoyo y fase&Generan el orden FL--RR--FR--RL y el
contacto previsto.\\
$k,T_s$&--, s&Índice y periodo discreto&Convierten ecuaciones continuas en un
algoritmo que se ejecuta a instantes definidos.\\
$\mat A_d,\mat B_d$&--&Dinámica discreta&Predicen el estado siguiente y permiten
analizar polos, controlabilidad y diseñar ganancias.\\
$\mat K$&según estado&Ganancia de realimentación&Transforma error de estado en
acción; en esta tesis el LQR actual es solo ejemplo nominal.\\
$\mat Q,\mat R$&pesos&Penalizaciones LQR&Expresan el compromiso entre error y
esfuerzo y hacen reproducible la sintonización.\\
$t_{LO},t_{TD}$&s&Instantes de despegue y aterrizaje&Permiten comprobar si la
pata físicamente ejecuta la fase publicada.\\
$\Delta t_{LO},\Delta t_{TD}$&s&Retardos de transición&Diagnostican dónde debe
corregirse la trayectoria.\\
$\gamma_i$&\%&Coincidencia de contacto&Cuantifica cumplimiento del plan sin
depender de apreciación visual.\\
$v_x$, roll, pitch&m/s, grados&Avance e inclinación&Resultados comparativos para
nominal frente a nominal más RL.\\
\bottomrule
\end{longtable}
\normalsize

\section{Parámetros nominales}
\begin{longtable}{p{4.6cm}p{1.7cm}p{2.8cm}p{5.3cm}}
\caption{Geometría, masa y entorno.}\label{tab:param}\\
\toprule Parámetro & Símbolo & Valor & Estado \\ \midrule\endfirsthead
\toprule Parámetro & Símbolo & Valor & Estado \\ \midrule\endhead
Separación longitudinal de caderas & $d_x$ & \SI{0.180}{m} & Referencia NovaSM3\\
Separación lateral de caderas & $d_y$ & \SI{0.120}{m} & Referencia NovaSM3\\
Longitud de coxa & $l_c$ & \SI{0.090}{m} & Referencia geométrica\\
Longitud de fémur & $l_f$ & \SI{0.105}{m} & Referencia geométrica\\
Longitud de tibia & $l_t$ & \SI{0.132}{m} & Referencia geométrica\\
Masa del cuerpo & $m_b$ & \SI{1.20}{kg} & Estimación nominal\\
Masa coxa/fémur/tibia/pie & --- & \SIlist{0.12;0.12;0.11;0.03}{kg} & Estimaciones nominales\\
Masa total & $m$ & \SI{2.72}{kg} & Referencia pública aproximada\\
Gravedad & $g$ & \SI{9.80665}{m.s^{-2}} & Convención estándar\\
Fricción pie--suelo & $\mu$ & $0.90$ & Parámetro de simulación\\ \bottomrule
\end{longtable}

\section{Transformaciones homogéneas}
\begin{tracebox}{T02 $\rightarrow$ A01 $\rightarrow$ C01 $\rightarrow$ E01}
La forma homogénea procede de robótica serial. Se eligió porque reúne rotación y
traslación sin perder el marco de referencia. En Nova se sustituyeron ejes REP--103,
posiciones de cadera y longitudes propias. La última columna se contrasta con
\texttt{forward\_leg} y con el URDF.
\end{tracebox}
Una transformación reúne orientación y traslación:
\begin{equation}
 {}^A\mat T_B=\begin{bmatrix}{}^A\mat R_B&{}^A\vect p_B\\
 0&0&0&1\end{bmatrix}. \label{eq:homogeneous}
\end{equation}
La cadena cuerpo--pie es
\begin{align}
 {}^B\mat T_{P_i}={}&{}^B\mat T_{H_i}\,
 \mat R_x(s_iq_{c,i})\mat T_y(s_il_c)\mat R_y(q_{f,i})
 \mat T_z(-l_f)\mat R_y(q_{t,i})\mat T_z(-l_t),\label{eq:chain}\\
 {}^W\vect p_i={}&{}^W\vect r_B+{}^W\mat R_B\,{}^B\vect p_i.\label{eq:worldfoot}
\end{align}
La última columna de \cref{eq:chain} debe coincidir numéricamente con la FK
trigonométrica. ROS 2 representa la orientación de base mediante cuaternión;
roll, pitch y yaw se derivan solo para métricas y visualización.

\section{Cinemática directa de una pata}
\begin{tracebox}{T01--T03 $\rightarrow$ A01 $\rightarrow$ C01 $\rightarrow$ E01}
Método: multiplicación de transformaciones y simplificación trigonométrica.
Objetivo de tesis: obtener la posición del pie a partir de las 12 posiciones
articulares para diseñar y verificar marchas cartesianas.
\end{tracebox}
Sea $a=s_iq_c$, $Q=q_f+q_t$ y
\begin{align}
 x &= -l_f\sin q_f-l_t\sin Q,\label{eq:fkx}\\
 z_p &= -l_f\cos q_f-l_t\cos Q,\label{eq:fkzp}\\
 y &= \cos(a)s_il_c-\sin(a)z_p,\label{eq:fky}\\
 z &= \sin(a)s_il_c+\cos(a)z_p.\label{eq:fkz}
\end{align}
Por tanto, $\vect p_i=[x,y,z]\trans=\FK_i(\vect q_i)$. Estas expresiones se
obtienen al formar primero la cadena planar fémur--tibia, sumar el
desplazamiento lateral de la coxa y rotar alrededor de $x$. El factor $s_i$
cambia simultáneamente lado y sentido de coxa, manteniendo una formulación
única para las cuatro patas.

\begin{workedexample}{cinemática directa de la pata delantera izquierda}
Se usa la postura nominal que realmente carga el controlador:
\[
 q_c=\SI{0.10}{rad},\quad q_f=\SI{0.42}{rad},\quad
 q_t=\SI{-0.84}{rad},\quad s_{FL}=+1.
\]
Primero, $Q=q_f+q_t=\SI{-0.42}{rad}$. Sustituyendo longitudes y ángulos:
\begin{align*}
 x&=-0.105\sin(0.42)-0.132\sin(-0.42)
    =\SI{0.0110095}{m},\\
 z_p&=-0.105\cos(0.42)-0.132\cos(-0.42)
    =\SI{-0.2164021}{m}.
\end{align*}
Como $a=sq_c=0.10$ rad y $s l_c=0.090$ m:
\begin{align*}
 y&=\cos(0.10)(0.090)-\sin(0.10)(-0.2164021)
   =\SI{0.1111545}{m},\\
 z&=\sin(0.10)(0.090)+\cos(0.10)(-0.2164021)
   =\SI{-0.2063360}{m}.
\end{align*}
Por tanto,
\[
 \boxed{{}^{H_{FL}}\vect p_{FL}=
 [0.0110095,\;0.1111545,\;-0.2063360]\trans\ \si{m}}.
\]
Este valor se obtiene con \texttt{forward\_leg} y constituye la posición
neutral usada para construir las trayectorias de gateo y marcha paso.
\end{workedexample}

\subsection{Centros de masa de la cadena}
Con $\vect e_z=[0,0,-1]\trans$:
\begin{align}
 \vect r_c &=\mat R_x(a)[0,s_il_c/2,0]\trans,\\
 \vect r_{ce}&=\mat R_x(a)[0,s_il_c,0]\trans,\\
 \vect u_f&=\mat R_x(a)\mat R_y(q_f)\vect e_z,
 &\vect u_t&=\mat R_x(a)\mat R_y(Q)\vect e_z,\\
 \vect r_f&=\vect r_{ce}+\tfrac{l_f}{2}\vect u_f,
 &\vect r_k&=\vect r_{ce}+l_f\vect u_f,\\
 \vect r_t&=\vect r_k+\tfrac{l_t}{2}\vect u_t,
 &\vect r_{pie}&=\vect r_k+l_t\vect u_t.
\end{align}

\section{Cinemática inversa, ramas y dominio}
\begin{tracebox}{T02--T03 $\rightarrow$ A02 $\rightarrow$ C02 $\rightarrow$ E01}
Método: eliminar primero la coxa y resolver después el triángulo fémur--tibia
por ley de cosenos. Se usa la rama de rodilla flexionada compatible con el URDF.
Objetivo: convertir cada punto cartesiano de la marcha en una referencia articular
alcanzable, continua y limitada.
\end{tracebox}
Para $\vect p=[x,y,z]\trans$, se elimina primero la coxa:
\begin{align}
 z_p&=-\sqrt{y^2+z^2-l_c^2},\label{eq:ikzp}\\
 a&=\operatorname{atan2}(z,y)-\operatorname{atan2}(z_p,s_il_c),
 &q_c&=s_i a.\label{eq:ikcoxa}
\end{align}
Si $y^2+z^2-l_c^2<0$, no existe solución real. Definiendo $d=-z_p$,
$b=-x$ y $r^2=b^2+d^2$:
\begin{align}
 D&=\frac{r^2-l_f^2-l_t^2}{2l_fl_t}, & -1\le D\le1,\label{eq:D}\\
 q_t&=-\arccos D,\label{eq:kneebranch}\\
 q_f&=\operatorname{atan2}(b,d)-
 \operatorname{atan2}\!\left(l_t\sin q_t,l_f+l_t\cos q_t\right).\label{eq:ikfemur}
\end{align}
La rama negativa representa rodilla flexionada. La rama positiva puede ser
matemáticamente válida, pero se descarta si invierte la rodilla, colisiona o
rompe continuidad. La solución se acepta si
\begin{equation}
 q_{j,\min}\le q_j\le q_{j,\max},\qquad
 \|\vect q_i(k)-\vect q_i(k-1)\|_\infty\le\delta_q.\label{eq:iklimits}
\end{equation}
La marcha nominal rechaza un objetivo inválido y conserva \texttt{stand}; no
lo satura silenciosamente. Una proyección futura deberá formularse como
$\vect p^*=\argmin_{\vect p\in\mathcal W_i}\|\vect p-\vect p_d\|$.

\begin{workedexample}{comprobación completa FK--IK}
Se usa el punto calculado anteriormente. Primero:
\begin{align*}
 z_p&=-\sqrt{0.1111545^2+(-0.2063360)^2-0.090^2}
      =\SI{-0.2164021}{m},\\
 a&=\operatorname{atan2}(-0.2063360,0.1111545)
 -\operatorname{atan2}(-0.2164021,0.090)=\SI{0.10}{rad},\\
 q_c&=s a=\SI{0.10}{rad}.
\end{align*}
Con $b=-0.0110095$ m y $d=0.2164021$ m:
\begin{align*}
 D&=\frac{b^2+d^2-0.105^2-0.132^2}{2(0.105)(0.132)}
 =0.6674628,\\
 q_t&=-\arccos(0.6674628)=\SI{-0.84}{rad},\\
 q_f&=\operatorname{atan2}(b,d)-\operatorname{atan2}(0.132\sin q_t,
 0.105+0.132\cos q_t)=\SI{0.42}{rad}.
\end{align*}
La IK devuelve $[0.10,0.42,-0.84]\trans$ rad. El error máximo calculado por
el script fue menor que $10^{-12}$ rad. Esto verifica ecuaciones y código; no
convierte las longitudes nominales en mediciones físicas.
\end{workedexample}

\section{Espacio de trabajo}
\begin{align}
 \mathcal W_i&=\{\vect p\in\R^3:\exists\vect q_i\in\mathcal Q_i,
 \FK_i(\vect q_i)=\vect p\},\label{eq:workspace}\\
 \mathcal W_i^{\mathrm{safe}}&=\{\FK_i(\vect q_i):
 \vect q_{\min}+\vect\epsilon_q\le\vect q_i\le
 \vect q_{\max}-\vect\epsilon_q,\ \sigma_{\min}(\mat J_i)\ge\epsilon_\sigma\}.
\end{align}
La construcción práctica muestrea el espacio articular, elimina colisiones y
márgenes insuficientes y genera una nube 3D. Deben incluirse cortes $x$--$z$ y
$y$--$z$, trayectoria nominal superpuesta y distancia mínima a límites.

\section{Jacobiano y singularidades}
\begin{tracebox}{T02,T04 $\rightarrow$ A03 $\rightarrow$ C03 $\rightarrow$ E02}
Método: derivación analítica de FK y principio de potencia virtual. Objetivo:
relacionar velocidades, detectar configuraciones mal condicionadas y convertir
fuerzas del pie en pares comparables con el MG996R.
\end{tracebox}
De \cref{eq:fkx,eq:fky,eq:fkz}:
\begin{equation}
 \dot{\vect p}_i=\mat J_i(\vect q_i)\dot{\vect q}_i,\qquad
 \mat J_i=\frac{\partial\vect p_i}{\partial\vect q_i}.\label{eq:jacdef}
\end{equation}
Sea $A=l_f\sin q_f+l_t\sin Q=-x$, $B=l_t\sin Q$ y $L=s_il_c$.
El Jacobiano analítico es
\begin{equation}
\mat J_i=\begin{bmatrix}
0&z_p&-l_t\cos Q\\
s_i[-\sin(a)L-\cos(a)z_p]&-\sin(a)A&-\sin(a)B\\
s_i[ \cos(a)L-\sin(a)z_p]& \cos(a)A& \cos(a)B
\end{bmatrix}.\label{eq:jacobian}
\end{equation}
La potencia virtual produce $\vect\tau_i=\mat J_i\trans\vect f_i$. La cercanía
a singularidad se evalúa mediante
\begin{equation}
 \kappa(\mat J_i)=\frac{\sigma_{\max}}{\sigma_{\min}},\qquad
 w_i=\sqrt{\det(\mat J_i\mat J_i\trans)}.\label{eq:singularity}
\end{equation}
Cuando $\sigma_{\min}<\epsilon_\sigma$ o $\kappa>\kappa_{\max}$ se reduce la
velocidad o se rechaza el objetivo. La pseudoinversa amortiguada es
\begin{equation}
\mat J_i^{\#}=\mat J_i\trans(\mat J_i\mat J_i\trans+\lambda^2\mat I)^{-1}.
\end{equation}

\begin{workedexample}{Jacobiano, singularidad y conversión fuerza--par}
Evaluando \cref{eq:jacobian} en la postura nominal FL se obtiene
\[
\mat J_{FL}=\begin{bmatrix}
0&-0.216402&-0.120528\\
0.206336&0.001099&0.005373\\
0.111155&-0.010955&-0.053555
\end{bmatrix}\ \si{m.rad^{-1}}.
\]
La descomposición en valores singulares calculada por NumPy entrega
\[
 (\sigma_1,\sigma_2,\sigma_3)=(0.252836,0.232233,0.038037)\ \si{m.rad^{-1}}.
\]
Por tanto,
\[
 \boxed{\kappa(\mat J)=\frac{0.252836}{0.038037}=6.647}.
\]
El valor es finito y la postura no es singular, aunque la dirección asociada a
$\sigma_3$ es menos manipulable. Si el pie ejerce
$\vect f=[1,0.5,5]\trans$ N, la potencia virtual da
\begin{align*}
 \vect\tau_f&=\mat J\trans\vect f\\
 &=\begin{bmatrix}
0&0.206336&0.111155\\-0.216402&0.001099&-0.010955\\
-0.120528&0.005373&-0.053555
\end{bmatrix}
\begin{bmatrix}1\\0.5\\5\end{bmatrix}\\
&=\boxed{[0.658941,-0.270625,-0.385618]\trans\ \si{N.m}}.
\end{align*}
Este término se resta en la dinámica inversa porque la fuerza externa ayuda o
se opone al par que debe generar el actuador según su signo.
\end{workedexample}

\subsection{Cinemática diferencial de segundo orden}
\begin{align}
 \ddot{\vect p}_i&=\mat J_i\ddot{\vect q}_i+\dot{\mat J}_i\dot{\vect q}_i,
 \label{eq:accel}\\
 \dot{\mat J}(\vect q,\dot{\vect q})&\approx
 \frac{\mat J(\vect q+h\dot{\vect q})-
 \mat J(\vect q-h\dot{\vect q})}{2h}.\label{eq:jdot}
\end{align}
$\dot{\mat J}$ debe verificarse contra diferencias finitas independientes. Se
comprueban límites articulares y cartesianos de posición, velocidad y
aceleración antes de emitir referencias.

\section{Modelo dinámico}
\begin{tracebox}{T04 $\rightarrow$ A04 $\rightarrow$ C04 $\rightarrow$ E03}
Método: energía cinética/potencial, Euler--Lagrange y símbolos de Christoffel.
La forma general se toma de dinámica multicuerpo; masas, inercias, signos y
geometría se adaptan a Nova. Objetivo: separar qué par se debe a aceleración,
velocidad, gravedad, fricción y contacto.
\end{tracebox}
La ecuación de base flotante es
\begin{equation}
 \mat M(\vect q)\ddot{\vect q}+\mat C(\vect q,\dot{\vect q})\dot{\vect q}
 +\vect g(\vect q)+\vect\tau_f
 =\mat S\trans\vect\tau+\mat J_c\trans\vect\lambda.\label{eq:dynamics}
\end{equation}
$\mat S$ selecciona las doce articulaciones porque la base no está actuada.
Para cada cuerpo $k$:
\begin{align}
 K&=\frac12\sum_k\left[m_k\vect v_{C_k}\trans\vect v_{C_k}+
 \vect\omega_k\trans{}^B\mat R_k\mat I_k{}^B\mat R_k\trans\vect\omega_k\right],
 \label{eq:kinetic}\\
 U&=\sum_km_kgz_{C_k},\label{eq:potential}\\
 M_{ab}&=\frac{\partial^2K}{\partial\dot q_a\partial\dot q_b},
 &g_a&=\frac{\partial U}{\partial q_a}.\label{eq:mg}
\end{align}
Para un paralelepípedo homogéneo de dimensiones $(a,b,c)$:
\begin{equation}
 \mat I_C=\frac{m}{12}\diag(b^2+c^2,a^2+c^2,a^2+b^2),\qquad
 \mat I_O=\mat I_C+m(\|\vect d\|^2\mat I-\vect d\vect d\trans).
\end{equation}
El símbolo de Christoffel de primera especie es
\begin{equation}
 c_{ijk}=\frac12\left(\frac{\partial M_{ij}}{\partial q_k}+
 \frac{\partial M_{ik}}{\partial q_j}-\frac{\partial M_{jk}}{\partial q_i}\right),
 \qquad [\vect c]_i=\sum_{j,k}c_{ijk}\dot q_j\dot q_k.\label{eq:christoffel}
\end{equation}
Las derivadas se evalúan mediante diferencias centrales con
$\varepsilon=10^{-5}$. La fricción pasiva suavizada es
\begin{equation}
 \vect\tau_f=\mat D\dot{\vect q}+\vect\tau_C\odot
 \tanh(\dot{\vect q}/v_s),\quad
 b=\SI{0.08}{N.m.s.rad^{-1}},\ \tau_C=\SI{0.03}{N.m},\ v_s=\SI{0.01}{rad.s^{-1}}.
\end{equation}
La dinámica inversa reducida resulta
\begin{equation}
 \vect\tau_i=\mat M_i\ddot{\vect q}_i+\vect c_i+\vect g_i+
 \vect\tau_{f,i}-\mat J_i\trans\vect f_i.\label{eq:inverse-dynamics}
\end{equation}
Se verifica $\mat M=\mat M\trans>0$ y
$\dot{\vect q}\trans(\dot{\mat M}-2\mat C)\dot{\vect q}=0$.

\begin{workedexample}{dinámica inversa de una pata, término por término}
Se emplean valores ejecutados por nuestro script reproducible:
\begin{align*}
 \vect q&=[0.10,0.42,-0.84]\trans\ \si{rad},\\
 \dot{\vect q}&=[0.10,-0.05,0.08]\trans\ \si{rad.s^{-1}},\\
 \ddot{\vect q}&=[0.20,0.10,-0.15]\trans\ \si{rad.s^{-2}},\\
 \vect f&=[1,0.5,5]\trans\ \si{N}.
\end{align*}
La matriz de masa obtenida de los Jacobianos de los centros de masa y las
inercias rotacionales es
\[
\mat M=\begin{bmatrix}
0.00904783&0.00035891&-0.00041176\\
0.00035891&0.00673928&0.00195619\\
-0.00041176&0.00195619&0.00316985
\end{bmatrix}\ \si{kg.m^2}.
\]
El primer término se calcula mediante multiplicación explícita:
\begin{align*}
 \mat M\ddot{\vect q}
 &=\mat M\begin{bmatrix}0.20\\0.10\\-0.15\end{bmatrix}
 =\begin{bmatrix}0.00190722\\0.00045228\\-0.00036221\end{bmatrix}\si{N.m}.
\end{align*}
Los restantes términos, calculados con las mismas funciones que pasan las
pruebas automatizadas, son
\begin{align*}
 \vect c&=[0.00002208,0.00000165,-0.00001087]\trans\ \si{N.m},\\
 \vect g&=[0.30982337,0.03891254,-0.04464199]\trans\ \si{N.m},\\
 \vect\tau_f&=[0.03800000,-0.03399728,0.03639999]\trans\ \si{N.m},\\
 \mat J\trans\vect f&=[0.65894065,-0.27062517,-0.38561842]\trans\ \si{N.m}.
\end{align*}
Sustituyendo sin agrupar términos:
\begin{align*}
 \vect\tau={}&
 \begin{bmatrix}0.00190722\\0.00045228\\-0.00036221\end{bmatrix}+
 \begin{bmatrix}0.00002208\\0.00000165\\-0.00001087\end{bmatrix}+
 \begin{bmatrix}0.30982337\\0.03891254\\-0.04464199\end{bmatrix}\\
 &+\begin{bmatrix}0.03800000\\-0.03399728\\0.03639999\end{bmatrix}
 -\begin{bmatrix}0.65894065\\-0.27062517\\-0.38561842\end{bmatrix}\\
 ={}&\boxed{[-0.309188,0.275994,0.377003]\trans\ \si{N.m}}.
\end{align*}
Los autovalores de $\mat M$ son
$(0.00226642,0.00758520,0.00910534)$ \si{kg.m^2}; todos son positivos. Así se
comprueba numéricamente que la energía cinética local es positiva para cualquier
velocidad articular distinta de cero.
\end{workedexample}

\section{Modelo del actuador MG996R}
\begin{table}[ht]\centering\caption{Datos de catálogo a \SI{6}{V}.}
\begin{tabular}{ll}\toprule Magnitud&Valor\\\midrule
Masa&\SI{55}{g}\\Par de bloqueo&\SI{1.0787}{N.m}\\
Velocidad sin carga&\SI{6.981}{rad.s^{-1}}\\
Corriente sin carga&\SI{0.170}{A}\\Corriente de bloqueo&\SI{1.40}{A}\\
Tensión&\SIrange{4.8}{6.6}{V}\\Banda muerta&\SI{1}{\micro s}\\\bottomrule
\end{tabular}\end{table}
Se distinguen articulación ideal, servo de posición y actuador físico:
\begin{align}
 q_{\mathrm{ref}}&=f_{\mathrm{cal}}(p_{\mathrm{PWM}};p_0,s,k,p_{\min},p_{\max}),\\
 \tau_{\mathrm{cmd}}&=\sat_{[-\tau_{\max},\tau_{\max}]}
 [K_p\operatorname{deadzone}(e,e_d)-K_d\dot q-\tau_f],\\
 |\tau_{\max}(\omega,V)|&=\tau_{\mathrm{stall}}(V)
 \max\!\left(0,1-\frac{|\omega|}{\omega_0(V)}\right).\label{eq:actuator}
\end{align}
Una aproximación inicial de corriente es
$I\approx I_0+(I_{\mathrm{stall}}-I_0)|\tau|/\tau_{\mathrm{stall}}$.
Centro, extremos, sentido, zona muerta, velocidad, corriente, temperatura y
respuesta bajo carga deben medirse por unidad. El par de bloqueo no es par
continuo seguro.

\begin{workedexample}{envolvente nominal de par y corriente}
Para la coxa del ejemplo, $|\omega|=\SI{0.10}{rad.s^{-1}}$. A \SI{6}{V},
$\tau_{stall}=\SI{1.0787315}{N.m}$ y
$\omega_0=\SI{6.981317}{rad.s^{-1}}$. Entonces
\begin{align*}
 \tau_{max}&=1.0787315\left(1-\frac{0.10}{6.981317}\right)
 =\SI{1.063280}{N.m}.
\end{align*}
El par requerido fue $|\tau_c|=\SI{0.309188}{N.m}$, equivalente al
$29.08\%$ de esa envolvente nominal. La corriente aproximada es
\begin{align*}
 I&=0.170+(1.40-0.170)\frac{0.309188}{1.0787315}\\
  &=\boxed{\SI{0.5225}{A}}.
\end{align*}
Para las tres articulaciones, el modelo estima
$[0.523,0.485,0.600]$ A. Estos valores sirven para dimensionamiento preliminar
y comparación, pero no sustituyen la medición con fuente limitada: el catálogo
no especifica par continuo, dispersión ni corriente transitoria de cada unidad.
\end{workedexample}

\section{Contacto pie--suelo}
\begin{tracebox}{T05,T11 $\rightarrow$ A06 $\rightarrow$ C06 $\rightarrow$ E06}
La complementariedad y Coulomb son la base física; el resorte--amortiguador es
una aproximación computable. Los sensores oficiales de Gazebo suministran la
evidencia observada. Objetivo: comprobar apoyo real y deslizamiento, no asumir
contacto porque la fase lo declara.
\end{tracebox}
La separación $\phi_i(\vect q)$ y las fuerzas satisfacen
\begin{align}
 \phi_i&\ge0,&\lambda_{n,i}&\ge0,&\phi_i\lambda_{n,i}&=0,\label{eq:complementarity}\\
 \sqrt{\lambda_{x,i}^2+\lambda_{y,i}^2}&\le\mu\lambda_{z,i}.\label{eq:frictioncone}
\end{align}
El modelo penalizado computable es
\begin{align}
 \delta_i&=\max(0,-\phi_i),\\
 \lambda_{n,i}&=\max(0,k_n\delta_i-c_n\dot\phi_i),\\
 \vect\lambda_{t,i}&=-\mu\lambda_{n,i}
 \tanh(\|\vect v_{t,i}\|/v_s)\frac{\vect v_{t,i}}
 {\|\vect v_{t,i}\|+\epsilon}.
\end{align}
Se usan nominalmente $k_n=\SI{12000}{N.m^{-1}}$,
$c_n=\SI{80}{N.s.m^{-1}}$, $\mu=0.90$ y $v_s=\SI{0.01}{m.s^{-1}}$.
Gazebo y MuJoCo resuelven contacto con algoritmos propios; no se afirma que
estos parámetros sean idénticos a los internos de cada solver.

\section{Centro de masa y estabilidad}
\begin{equation}
 {}^B\vect r_{COM}=\frac{m_b\vect r_b+
 \sum_{i=1}^4\sum_jm_{ij}({}^B\vect p_{H_i}+\vect r_{ij})}{m_{\mathrm{total}}}.
\end{equation}
Para contactos ordenados antihorariamente, la distancia firmada a la arista
$\vect e_j=\vect p_{j+1}-\vect p_j$ es
\begin{equation}
 d_j=\frac{e_{x,j}(c_y-p_{y,j})-e_{y,j}(c_x-p_{x,j})}{\|\vect e_j\|},
 \qquad m_s=\min_jd_j.\label{eq:staticmargin}
\end{equation}
$m_s>0$ indica COM dentro del polígono. Se requieren tres contactos no
colineales. Para movimientos rápidos se usa balance centroidal:
\begin{align}
 m\ddot{\vect r}_{COM}&=\sum_i\vect f_i+m\vect g,\\
 \dot{\vect L}_{COM}&=\sum_i(\vect r_i-\vect r_{COM})\times\vect f_i.
\end{align}
El galope no se valida mediante margen estático; requerirá dinámica centroidal,
ZMP u otro criterio justificado.

\begin{workedexample}{centro de masa y margen con cuatro apoyos}
Para las cuatro patas en $[0.10,0.42,-0.84]\trans$ rad, la suma ponderada de
los centros de masa nominales produce
\[
 {}^B\vect r_{COM}=\boxed{[-0.0058646,0,-0.0388199]\trans\ \si{m}}.
\]
Al sumar la posición de cada cadera y la FK local, los contactos proyectados son
\begin{align*}
 FL&=(0.101010,\phantom{-}0.171155),&
 FR&=(0.101010,-0.171155),\\
 RL&=(-0.078990,\phantom{-}0.171155),&
 RR&=(-0.078990,-0.171155)\quad\si{m}.
\end{align*}
El COM proyectado está a \SI{0.0731259}{m} de la arista trasera, a
\SI{0.0958646}{m} de la delantera y a \SI{0.1711545}{m} de cada arista lateral.
El mínimo es
\[
 \boxed{m_s=\SI{0.0731259}{m}>0}.
\]
Así se obtuvo el margen nominal reportado. Es válido para cuatro contactos
ideales en esa postura; durante el gateo debe recalcularse con los contactos
medidos y no con los tres apoyos que el plan simplemente supone.
\end{workedexample}

\section{Generación matemática de marchas}
\begin{tracebox}{T02,T03 $\rightarrow$ A07/A08 $\rightarrow$ C07/C08 $\rightarrow$ E07/E08}
Método: parametrización por fase, apoyo lineal, oscilación parabólica e IK por
muestra. Objetivo: producir dos líneas base deterministas y repetibles contra
las cuales se comparará la corrección RL.
\end{tracebox}
La fase de cada pata y el indicador ideal de contacto son
\begin{equation}
 \varphi_i(t)=\operatorname{mod}\left(\frac{t}{T}+\Delta_i,1\right),\qquad
 c_i(\varphi)=\begin{cases}1,&0\le\varphi<\beta,\\0,&\beta\le\varphi<1.
 \end{cases}\label{eq:gaitphase}
\end{equation}
Para gateo, $\beta=0.75$ y
$(\Delta_{FL},\Delta_{RR},\Delta_{FR},\Delta_{RL})=(0,0.25,0.50,0.75)$.
Sea $L$ la longitud, $H$ la elevación, $z_0$ el suelo y $u$ fase local:
\begin{align}
 0\le\varphi<\beta:\quad &u=\varphi/\beta,&
 x&=x_0+L(1/2-u),&z&=z_0,\label{eq:stance}\\
 \beta\le\varphi<1:\quad &u=(\varphi-\beta)/(1-\beta),&
 x&=x_0+L(u-1/2),&z&=z_0+4Hu(1-u).\label{eq:swing}
\end{align}
La línea base usa $N=24$, $\Delta t=\SI{0.18}{s}$, $L=\SI{0.018}{m}$ y
$H=\SI{0.014}{m}$, por lo que
\begin{equation}T=N\Delta t=\SI{4.32}{s}.\end{equation}
La cadena completa es
\begin{equation}
 (L,H,T,\beta,\Delta_i)\rightarrow\vect p_i^d(t)\rightarrow\IK
 \rightarrow\vect q_i^d(t)\rightarrow\text{controlador}\rightarrow
 \vect\tau_i\rightarrow\text{movimiento/contacto}.\label{eq:motionchain}
\end{equation}

\subsection{Trayectorias suaves}
La parábola de \cref{eq:swing} es útil, pero no garantiza aceleración continua
en todos los cambios. Una interpolación quíntica es
\begin{align}
 b(u)&=10u^3-15u^4+6u^5,\\
 \vect p(t)&=\vect p_0+b(t/T)(\vect p_1-\vect p_0),\label{eq:quintic}\\
 \vect p^-&=\vect p^+,\quad\dot{\vect p}^-=\dot{\vect p}^+,\quad
 \ddot{\vect p}^-=\ddot{\vect p}^+.
\end{align}
Todo cambio de $L,H,T,\beta,N$ exige repetir alcanzabilidad, singularidad,
continuidad, límites, estabilidad y ensayo mínimo de diez ciclos.

\begin{workedexample}{una oscilación discreta de la pata FL}
Con $N=24$ y $\beta=0.75$, cada pata dispone de
$(1-\beta)N=6$ referencias de oscilación. Para FL, las fases discretas usan
$u\in\{0,1/6,2/6,3/6,4/6,5/6\}$. La altura relativa sigue
$h(u)=4Hu(1-u)$, con $H=0.014$ m:
\begin{center}
\begin{tabular}{c@{\qquad}c@{\qquad}c}
\toprule Muestra local&$u$&$h(u)$ [mm]\\\midrule
0&0&0.000\\1&1/6&7.778\\2&2/6&12.444\\
3&3/6&14.000\\4&4/6&12.444\\5&5/6&7.778\\\bottomrule
\end{tabular}
\end{center}
El código alcanza exactamente \SI{14}{mm}, pero la última referencia todavía
está \SI{7.778}{mm} elevada y la siguiente muestra vuelve inmediatamente a
apoyo. Esta discontinuidad discreta ayuda a explicar por qué la etiqueta de
aterrizaje no coincide con el contacto físico. No se oculta: es una propiedad
de la trayectoria que debe corregirse en la siguiente versión.
\end{workedexample}

\begin{projectresult}{cadencia comprobada en Gazebo}
El cálculo nominal es $24(0.18)=\SI{4.32}{s}$. La línea base corregida midió
\SI{4.320013}{s} por ciclo. El error absoluto medio fue aproximadamente
\[
 |4.320013-4.320000|=\SI{13}{\micro s},
\]
equivalente a $0.00030\%$. La repetición midió \SI{4.320002}{s} en el ensayo de
contactos, confirmando que el cambio de deadline eliminó la deriva acumulativa.
\end{projectresult}

\section{Fundamentos del control discreto}
\label{sec:discrete-control}
\begin{tracebox}{T07,T10 $\rightarrow$ A09 $\rightarrow$ C09 $\rightarrow$ E09}
La teoría de sistemas muestreados aporta $k$, $T_s$, ZOH y plano Z. La
adaptación distingue temporizador, referencia y controlador articular ROS 2.
Objetivo: demostrar que la cadencia implementada corresponde a un algoritmo
discreto verificable y preparar futuros lazos físicos.
\end{tracebox}
Esta sección desarrolla el sistema de control como se estudiaría en un curso de
control discreto y lo conecta con la implementación del robot. Es esencial
distinguir tres escalas de tiempo:
\begin{enumerate}
 \item el temporizador de software consulta el reloj cada $T_{timer}=\SI{0.02}{s}$;
 \item el secuenciador publica una nueva referencia cada $T_r=\SI{0.18}{s}$;
 \item el controlador articular interno del simulador opera con su propio
 periodo $T_c$, normalmente menor que $T_r$.
\end{enumerate}
El temporizador de \SI{20}{ms} no convierte la marcha en un controlador de
\SI{50}{Hz}: solo comprueba si venció la próxima fase. La señal de referencia
nominal cambia a $1/T_r=\SI{5.556}{Hz}$. El instante programado se actualiza
desde el vencimiento anterior,
\begin{equation}
 t_{k+1}^{d}=t_k^{d}+T_r,
 \qquad e_k^t=t_k^{pub}-t_k^d,
 \label{eq:deadline}
\end{equation}
para evitar que el jitter $e_k^t$ se acumule ciclo tras ciclo.

\subsection{De tiempo continuo a secuencias}
Sea $x(t)$ una señal continua. Un muestreador ideal produce
\begin{equation}
 x[k]=x(kT_s),\qquad k=0,1,2,\ldots,
 \label{eq:sampling}
\end{equation}
donde $T_s$ es el periodo y $f_s=1/T_s$ la frecuencia de muestreo. La señal
muestreada puede representarse mediante un tren de impulsos
\begin{equation}
 x_s(t)=x(t)\sum_{k=-\infty}^{\infty}\delta(t-kT_s).
\end{equation}
Su espectro contiene réplicas separadas por $\omega_s=2\pi/T_s$. Para evitar
aliasing, una señal limitada a $\omega_b$ debe satisfacer idealmente
\begin{equation}
 \omega_s>2\omega_b.
 \label{eq:nyquist}
\end{equation}
En un robot real se requiere además un filtro analógico antialias antes del
conversor. El cumplimiento de Nyquist es necesario, pero no garantiza buen
control: en práctica se selecciona una frecuencia de muestreo entre 10 y 20
veces el ancho de banda de lazo cerrado, sujeta a cómputo, ruido y retardos.

Entre muestras, un retenedor de orden cero (ZOH) mantiene constante el comando:
\begin{equation}
 u(t)=u[k],\qquad kT_s\le t<(k+1)T_s.
 \label{eq:zoh}
\end{equation}
Esta hipótesis representa adecuadamente una salida digital de posición o PWM
actualizada periódicamente. Las referencias ROS 2 se interpolan internamente,
por lo que el comportamiento exacto debe identificarse como ZOH, rampa u otra
interpolación antes de diseñar un controlador físico equivalente.

\subsection{Ecuaciones en diferencias}
Un sistema discreto lineal causal puede expresarse como
\begin{equation}
 y[k]+a_1y[k-1]+\cdots+a_ny[k-n]
 =b_0u[k]+b_1u[k-1]+\cdots+b_mu[k-m].
 \label{eq:difference}
\end{equation}
Esta ecuación indica exactamente qué memoria necesita una implementación. Por
ejemplo, el filtro de primer orden
\begin{equation}
 y[k]=\alpha y[k-1]+(1-\alpha)u[k],\qquad 0<\alpha<1,
\end{equation}
almacena una salida anterior. Si se parte de una constante de tiempo $\tau$ y
se usa discretización exacta, $\alpha=e^{-T_s/\tau}$.

\subsection{Transformada Z}
La transformada Z unilateral de una secuencia es
\begin{equation}
 X(z)=\mathcal Z\{x[k]\}=\sum_{k=0}^{\infty}x[k]z^{-k}.
 \label{eq:ztransform}
\end{equation}
El operador $z^{-1}$ representa un retardo de una muestra. Con condiciones
iniciales nulas,
\begin{equation}
 \mathcal Z\{x[k-r]\}=z^{-r}X(z).
\end{equation}
Aplicar la transformada a \cref{eq:difference} produce la función de
transferencia de pulsos
\begin{equation}
 G(z)=\frac{Y(z)}{U(z)}=
 \frac{b_0+b_1z^{-1}+\cdots+b_mz^{-m}}
 {1+a_1z^{-1}+\cdots+a_nz^{-n}}.
 \label{eq:pulse-transfer}
\end{equation}
Los polos son las raíces del denominador y los ceros las del numerador. La
región de convergencia debe incluir el círculo unitario para que una respuesta
causal racional sea BIBO estable.

\subsection{Correspondencia entre los planos $s$ y $z$}
El muestreo relaciona los modos continuos y discretos mediante
\begin{equation}
 z=e^{sT_s}.
 \label{eq:szmap}
\end{equation}
Si $s=\sigma+j\omega$, entonces $|z|=e^{\sigma T_s}$ y
$\angle z=\omega T_s$. Por tanto, el semiplano continuo estable $\sigma<0$ se
mapea al interior del círculo unitario $|z|<1$; el eje imaginario se mapea al
círculo unitario y el semiplano inestable queda fuera.

Para recuperar especificaciones aproximadas de un polo discreto
$z_p=r e^{j\theta}$:
\begin{equation}
 s_p=\frac{\ln z_p}{T_s}
 =\frac{\ln r}{T_s}+j\frac{\theta}{T_s}.
\end{equation}
De $s_p=-\zeta\omega_n+j\omega_n\sqrt{1-\zeta^2}$ se obtienen amortiguamiento,
frecuencia natural, sobreimpulso y tiempo de establecimiento esperados.

\section{Discretización de modelos del robot}
\subsection{Discretización exacta con ZOH}
Para el modelo continuo linealizado
\begin{equation}
 \dot{\vect x}(t)=\mat A\vect x(t)+\mat B\vect u(t),\qquad
 \vect y(t)=\mat C\vect x(t)+\mat D\vect u(t),
 \label{eq:sscontinuous}
\end{equation}
y entrada constante durante una muestra, la solución exacta es
\begin{align}
 \vect x[k+1]&=\mat A_d\vect x[k]+\mat B_d\vect u[k],\label{eq:ssdiscrete}\\
 \mat A_d&=e^{\mat A T_s},&
 \mat B_d&=\int_0^{T_s}e^{\mat A\tau}\mat B\,d\tau,\label{eq:exactdisc}\\
 \vect y[k]&=\mat C_d\vect x[k]+\mat D_d\vect u[k],
\end{align}
con $\mat C_d=\mat C$ y $\mat D_d=\mat D$. No debe calcularse
$\mat A^{-1}(\mat A_d-\mat I)\mat B$ si $\mat A$ es singular. Una forma
numéricamente robusta usa la exponencial por bloques:
\begin{equation}
 \exp\!\left(
 \begin{bmatrix}\mat A&\mat B\\\mat 0&\mat 0\end{bmatrix}T_s
 \right)=
 \begin{bmatrix}\mat A_d&\mat B_d\\\mat 0&\mat I\end{bmatrix}.
 \label{eq:blockexp}
\end{equation}

\subsection{Euler hacia delante, Euler hacia atrás y Tustin}
Cuando se requiere una aproximación explícita:
\begin{align}
 \text{Euler directo: }&s\approx\frac{z-1}{T_s},&
 \mat A_d&\approx\mat I+T_s\mat A,&\mat B_d&\approx T_s\mat B;\\
 \text{Euler inverso: }&s\approx\frac{z-1}{T_sz},\\
 \text{Tustin: }&s\approx\frac{2}{T_s}\frac{z-1}{z+1}.
\end{align}
Euler directo puede volver inestable una simulación de un sistema estable si
$T_s$ es grande. Euler inverso es estable para una clase mayor de problemas,
pero introduce amortiguamiento numérico. Tustin mapea el semiplano izquierdo al
interior del círculo unitario y suele ser apropiado para discretizar filtros y
controladores. Si interesa preservar una frecuencia $\omega_p$, se aplica
prewarping sustituyendo $2/T_s$ por $\omega_p/\tan(\omega_pT_s/2)$.

\subsection{Linealización articular alrededor de una postura}
Partiendo de la dinámica reducida
\begin{equation}
 \mat M(\vect q)\ddot{\vect q}+\vect c(\vect q,\dot{\vect q})
 +\vect g(\vect q)=\vect\tau+\mat J\trans\vect f,
\end{equation}
se elige un equilibrio $(\vect q_0,\dot{\vect q}_0=0,\vect\tau_0)$ y se
definen perturbaciones $\delta\vect q=\vect q-\vect q_0$. Al primer orden,
\begin{equation}
 \mat M_0\delta\ddot{\vect q}+\mat D_0\delta\dot{\vect q}
 +\mat K_0\delta\vect q=\delta\vect\tau+\mat J_0\trans\delta\vect f,
 \label{eq:linearjoint}
\end{equation}
donde $\mat K_0=\partial\vect g/\partial\vect q|_0$ y $\mat D_0$ agrupa
amortiguamiento físico y del actuador. Con
$\vect x=[\delta\vect q\trans,\delta\dot{\vect q}\trans]\trans$:
\begin{equation}
 \mat A=\begin{bmatrix}\mat 0&\mat I\\-\mat M_0^{-1}\mat K_0&
 -\mat M_0^{-1}\mat D_0\end{bmatrix},\quad
 \mat B=\begin{bmatrix}\mat 0\\\mat M_0^{-1}\end{bmatrix}.
 \label{eq:jointss}
\end{equation}
Este modelo permite diseñar control local, pero cambia con la postura y el
contacto. Deben estudiarse varios puntos de operación o emplearse programación
de ganancias; una sola linealización no representa el ciclo completo.

\begin{workedexample}{linealización y discretización del eje de fémur Nova}
Para mostrar el procedimiento con nuestros parámetros, se aísla localmente el
eje de fémur alrededor de la postura nominal. Se usan
\begin{align*}
 M_f&=M_{22}=\SI{0.006739276}{kg.m^2},\\
 D_f&=\SI{0.08}{N.m.s.rad^{-1}},\\
 K_g&=\left.\frac{\partial g_f}{\partial q_f}\right|_{\vect q_0}
 \approx\frac{g_f(q_f+10^{-6})-g_f(q_f-10^{-6})}{2(10^{-6})}
 =\SI{0.287048}{N.m.rad^{-1}}.
\end{align*}
La ecuación perturbada es
\[
 M_f\delta\ddot q_f+D_f\delta\dot q_f+K_g\delta q_f=\delta\tau_f.
\]
Dividiendo cada término por $M_f$:
\[
 \delta\ddot q_f=-42.593324\,\delta q_f
 -11.870711\,\delta\dot q_f+148.383894\,\delta\tau_f.
\]
Con $\vect x=[\delta q_f,\delta\dot q_f]\trans$:
\[
 \mat A=\begin{bmatrix}0&1\\-42.593324&-11.870711\end{bmatrix},\qquad
 \mat B=\begin{bmatrix}0\\148.383894\end{bmatrix}.
\]
Se elige $T_s=\SI{0.02}{s}$, correspondiente al periodo de consulta del nodo,
y se evalúa la exponencial por bloques de \cref{eq:blockexp}. El resultado es
\[
 \boxed{\mat A_d=\begin{bmatrix}
0.992128&0.017753\\-0.756144&0.781392
\end{bmatrix}},\qquad
 \boxed{\mat B_d=\begin{bmatrix}0.027423\\2.634206\end{bmatrix}}.
\]
Los polos abiertos son $0.886760\pm j0.048178$, cuyo módulo es $0.888068<1$.
La matriz de controlabilidad
$\mathcal C=[\mat B_d,\mat A_d\mat B_d]$ tiene rango 2. El modo local es, por
tanto, estable y controlable bajo estas hipótesis nominales.
\end{workedexample}

\section{Análisis de sistemas discretos}
\subsection{Respuesta libre, forzada y estabilidad}
La solución de \cref{eq:ssdiscrete} es
\begin{equation}
 \vect x[k]=\mat A_d^k\vect x[0]
 +\sum_{j=0}^{k-1}\mat A_d^{k-1-j}\mat B_d\vect u[j].
 \label{eq:discretesolution}
\end{equation}
El primer término es la respuesta libre y el segundo la respuesta forzada. El
sistema lineal autónomo es asintóticamente estable si y solo si
\begin{equation}
 \rho(\mat A_d)=\max_i|\lambda_i(\mat A_d)|<1.
 \label{eq:schur}
\end{equation}
Para sistemas racionales escalares, todos los polos deben quedar estrictamente
dentro del círculo unitario, sin cancelaciones inestables ocultas.

\subsection{Criterio de Jury}
El criterio de Jury comprueba estabilidad sin calcular explícitamente las
raíces. Para un polinomio de segundo orden
$P(z)=z^2+a_1z+a_2$, las condiciones necesarias y suficientes son
\begin{equation}
 |a_2|<1,\qquad 1+a_1+a_2>0,\qquad 1-a_1+a_2>0.
 \label{eq:jury2}
\end{equation}
En órdenes mayores se construye la tabla de Jury. En este proyecto se prefiere
verificar numéricamente polos y margen de robustez, conservando Jury como una
comprobación analítica para modelos de bajo orden.

\subsection{Estabilidad de Lyapunov}
Una matriz $\mat A_d$ es Schur estable si, para toda $\mat Q=\mat Q\trans>0$,
existe una única $\mat P=\mat P\trans>0$ que satisface
\begin{equation}
 \mat A_d\trans\mat P\mat A_d-\mat P=-\mat Q.
 \label{eq:dlyap}
\end{equation}
Con $V[k]=\vect x[k]\trans\mat P\vect x[k]$ se obtiene
$\Delta V=-\vect x[k]\trans\mat Q\vect x[k]<0$. Para el robot híbrido, una
misma $\mat P$ para todos los modos aportaría una garantía común fuerte; si se
usan funciones distintas por contacto, deben verificarse las transiciones.

\subsection{Error permanente y tipo del sistema}
Para realimentación unitaria, $E(z)=R(z)/(1+L(z))$. Si los polos de
$(z-1)E(z)$ quedan dentro del círculo unitario, el teorema del valor final da
\begin{equation}
 e_{ss}=\lim_{k\to\infty}e[k]=\lim_{z\to1}(z-1)E(z).
 \label{eq:finalvalue}
\end{equation}
El número de polos de lazo abierto en $z=1$ determina el tipo discreto y el
error ante escalón, rampa o parábola. Añadir acción integral incorpora un polo
en $z=1$, pero aumenta el orden y exige rediseñar estabilidad y anti-windup.

\subsection{Controlabilidad y observabilidad}
Para $n$ estados, las matrices son
\begin{align}
 \mathcal C&=[\mat B_d,\mat A_d\mat B_d,\ldots,
 \mat A_d^{n-1}\mat B_d],& \rank(\mathcal C)&=n,\label{eq:ctrb}\\
 \mathcal O&=\begin{bmatrix}\mat C_d\\\mat C_d\mat A_d\\\vdots\\
 \mat C_d\mat A_d^{n-1}\end{bmatrix},&\rank(\mathcal O)&=n.
 \label{eq:obsv}
\end{align}
Controlabilidad significa que una secuencia finita puede mover el estado;
observabilidad, que las salidas permiten reconstruirlo. Rango numérico debe
evaluarse con valores singulares y condición, no solo con igualdad exacta. Un
sistema casi no controlable u observable será sensible a ruido y redondeo.

\section{Diseño de controladores digitales}
\subsection{PID discreto}
Con error $e[k]=r[k]-y[k]$, una realización posicional sencilla es
\begin{align}
 I[k]&=I[k-1]+T_s e[k],\\
 D_f[k]&=\alpha D_f[k-1]+(1-\alpha)\frac{e[k]-e[k-1]}{T_s},\\
 u_0[k]&=K_pe[k]+K_iI[k]+K_dD_f[k],\\
 u[k]&=\sat_{[u_{min},u_{max}]}(u_0[k]).
 \label{eq:pid}
\end{align}
La derivada debe filtrarse y, para evitar golpes ante cambios de consigna,
puede aplicarse a la medida en lugar del error. Una realización incremental,
útil para actualizar directamente el comando, es
\begin{equation}
 \Delta u[k]=K_p(e[k]-e[k-1])+K_iT_se[k]
 +\frac{K_d}{T_s}(e[k]-2e[k-1]+e[k-2]).
\end{equation}
Los coeficientes dependen de $T_s$: cambiar la frecuencia sin recalcularlos
cambia el controlador físico.

\subsection{Saturación y anti-windup}
Los MG996R y las referencias articulares tienen límites de posición, velocidad,
par y PWM. Si el integrador continúa creciendo durante saturación aparece
windup. La realimentación por seguimiento corrige
\begin{equation}
 I[k]=I[k-1]+T_s\left(e[k]+K_{aw}(u[k-1]-u_0[k-1])\right).
 \label{eq:antiwindup}
\end{equation}
También puede usarse integración condicional. La saturación debe aparecer en
simulación, entrenamiento y supervisor; no basta con recortar PWM al final.

\subsection{Realimentación de estados y asignación de polos}
Para el sistema controlable se propone
\begin{equation}
 \vect u[k]=-\mat K\vect x[k]+\mat N_r\vect r[k].
 \label{eq:statefeedback}
\end{equation}
La dinámica cerrada es $\mat A_d-\mat B_d\mat K$. La ganancia $\mat K$ puede
obtenerse por asignación de polos. Para una referencia constante y salida
cuadrada compatible, el precompensador satisface
\begin{equation}
 \mat N_r=\left[\mat C_d(\mat I-\mat A_d+\mat B_d\mat K)^{-1}\mat B_d
 +\mat D_d\right]^{-1}.
\end{equation}
No se deben escoger polos arbitrariamente rápidos: elevan esfuerzo, amplifican
ruido y hacen dominante el retardo computacional.

\subsection{Acción integral en espacio de estados}
Para eliminar error constante se define
\begin{equation}
 \vect\xi[k+1]=\vect\xi[k]+T_s(\vect r[k]-\vect y[k]).
\end{equation}
El sistema aumentado es
\begin{equation}
 \begin{bmatrix}\vect x[k+1]\\\vect\xi[k+1]\end{bmatrix}=
 \begin{bmatrix}\mat A_d&\mat 0\\-T_s\mat C_d&\mat I\end{bmatrix}
 \begin{bmatrix}\vect x[k]\\\vect\xi[k]\end{bmatrix}+
 \begin{bmatrix}\mat B_d\\-T_s\mat D_d\end{bmatrix}\vect u[k]
 +\begin{bmatrix}\mat 0\\T_s\mat I\end{bmatrix}\vect r[k].
 \label{eq:augmented}
\end{equation}
Se diseña $\vect u=-\mat K_x\vect x-\mat K_i\vect\xi$ verificando la
controlabilidad aumentada y usando anti-windup.

\subsection{Regulador cuadrático lineal discreto}
El LQR minimiza
\begin{equation}
 J=\sum_{k=0}^{\infty}
 \left(\vect x[k]\trans\mat Q\vect x[k]
 +\vect u[k]\trans\mat R\vect u[k]\right),
 \label{eq:lqrcost}
\end{equation}
con $\mat Q\succeq0$ y $\mat R\succ0$. La ecuación algebraica de Riccati es
\begin{equation}
 \mat P=\mat A_d\trans\mat P\mat A_d-
 \mat A_d\trans\mat P\mat B_d(\mat R+\mat B_d\trans\mat P\mat B_d)^{-1}
 \mat B_d\trans\mat P\mat A_d+\mat Q,
\end{equation}
y
\begin{equation}
 \mat K=(\mat R+\mat B_d\trans\mat P\mat B_d)^{-1}
 \mat B_d\trans\mat P\mat A_d.
\end{equation}
$\mat Q$ y $\mat R$ no son ``ganancias mágicas'': codifican el compromiso
entre error y esfuerzo y deben normalizarse según rangos físicos.

\begin{workedexample}{diseño LQR discreto sobre el fémur nominal}
Se continúa el ejemplo anterior. Para penalizar con mayor fuerza el error
angular que la velocidad se selecciona, únicamente como diseño reproducible,
\[
 \mat Q=\diag(400,4),\qquad \mat R=[0.5].
\]
El procedimiento ejecutado por SciPy es:
\begin{enumerate}
 \item resolver la ecuación de Riccati discreta para $\mat P$;
 \item calcular
 $\mat K=(\mat R+\mat B_d\trans\mat P\mat B_d)^{-1}
 \mat B_d\trans\mat P\mat A_d$;
 \item formar $\mat A_{cl}=\mat A_d-\mat B_d\mat K$;
 \item verificar que todos los polos cerrados tengan módulo menor que uno.
\end{enumerate}
El resultado numérico fue
\[
 \boxed{\mat K=[3.115982\quad0.325030]},
 \qquad\lambda(\mat A_{cl})=\boxed{\{0.817918,0.013957\}}.
\]
Para una perturbación inicial de $10^\circ=0.174533$ rad y velocidad cero, el
primer comando es
\[
 \tau[0]=-\mat K\vect x[0]
 =-3.115982(0.174533)=\SI{-0.54384}{N.m}.
\]
Al iterar $\vect x[k+1]=\mat A_{cl}\vect x[k]$, el ángulo entra y permanece
dentro de $\pm0.2^\circ$ a los \SI{0.40}{s}. El par máximo es
\SI{0.54384}{N.m}, cerca del 50\% del par de bloqueo nominal. Esto demuestra el
proceso de diseño, pero no autoriza esa ganancia en hardware: faltan identificar
el servo cerrado real, retardos, saturación continua y flexibilidad.
\end{workedexample}

\section{Estimación de estado}
\subsection{Observador de Luenberger}
Si no se miden todas las velocidades o estados corporales,
\begin{equation}
 \hat{\vect x}[k+1]=\mat A_d\hat{\vect x}[k]+\mat B_d\vect u[k]
 +\mat L(\vect y[k]-\hat{\vect y}[k]),
 \quad \hat{\vect y}[k]=\mat C_d\hat{\vect x}[k]+\mat D_d\vect u[k].
 \label{eq:observer}
\end{equation}
El error evoluciona con $\mat A_d-\mat L\mat C_d$. Sus polos suelen elegirse
más rápidos que los del control, pero no tanto que amplifiquen cuantización,
impactos y ruido de IMU.

\subsection{Filtro de Kalman discreto}
Con ruidos $\vect w[k]\sim(0,\mat Q_w)$ y
$\vect v[k]\sim(0,\mat R_v)$:
\begin{align}
 \hat{\vect x}_{k|k-1}&=\mat A_d\hat{\vect x}_{k-1|k-1}+\mat B_d\vect u[k-1],\\
 \mat P_{k|k-1}&=\mat A_d\mat P_{k-1|k-1}\mat A_d\trans+\mat Q_w,\\
 \mat K_k&=\mat P_{k|k-1}\mat C_d\trans
 (\mat C_d\mat P_{k|k-1}\mat C_d\trans+\mat R_v)^{-1},\\
 \hat{\vect x}_{k|k}&=\hat{\vect x}_{k|k-1}+\mat K_k
 (\vect y[k]-\mat C_d\hat{\vect x}_{k|k-1}),\\
 \mat P_{k|k}&=(\mat I-\mat K_k\mat C_d)\mat P_{k|k-1}.
\end{align}
Las covarianzas se estiman con datos en reposo y movimiento; no deben ajustarse
solo para producir curvas visualmente suaves. Para el modelo no lineal pueden
evaluarse EKF o UKF, conservando pruebas de consistencia de innovación.

\section{Diseño en el dominio Z y respuesta en frecuencia}
\subsection{Lugar geométrico de las raíces}
Para $L(z)=K G(z)H(z)$, la ecuación característica es
\begin{equation}
 1+K G(z)H(z)=0.
\end{equation}
Un punto $z_0$ pertenece al lugar de raíces si cumple
\begin{align}
 \angle G(z_0)H(z_0)&=(2\ell+1)\pi,\\
 K&=\frac{1}{|G(z_0)H(z_0)|},\qquad K>0.
\end{align}
Las ramas parten de los polos de lazo abierto y terminan en sus ceros, incluidos
los ceros en infinito. A diferencia del plano $s$, la frontera de estabilidad
es el círculo unitario. Se deben superponer líneas de amortiguamiento y
frecuencia obtenidas mediante $s=\ln(z)/T_s$, no copiarlas directamente del
plano continuo.

Un compensador discreto de adelanto o atraso puede escribirse
\begin{equation}
 C(z)=K_c\frac{z-z_0}{z-p_0}.
\end{equation}
En adelanto, el par polo--cero aporta fase para mover polos dominantes; en
atraso mejora ganancia de baja frecuencia y error permanente. Todo cero o polo
añadido modifica respuesta transitoria, sensibilidad al ruido y realización
numérica, por lo que se verifica el controlador completo.

\subsection{Respuesta en frecuencia discreta}
La respuesta se evalúa sobre el círculo unitario:
\begin{equation}
 G(e^{j\Omega})=G(z)|_{z=e^{j\Omega}},\qquad
 \Omega=\omega T_s\in[0,\pi].
\end{equation}
$\Omega=\pi$ corresponde a Nyquist. El lazo
$L(e^{j\Omega})$ permite obtener margen de ganancia, margen de fase y ancho de
banda. Para un modelo SISO:
\begin{equation}
 S(z)=\frac{1}{1+L(z)},\qquad
 T(z)=\frac{L(z)}{1+L(z)},\qquad S(z)+T(z)=1.
\end{equation}
$S$ describe rechazo de perturbaciones e incertidumbre a baja frecuencia;
$T$ describe seguimiento y transmisión de ruido/modelo a alta frecuencia. No
es posible hacer ambas arbitrariamente pequeñas en todas las frecuencias.

\subsection{Criterio de Nyquist discreto}
El criterio conserva la lógica de rodeos del punto $-1$, pero el contorno se
construye a partir del círculo unitario y considera polos fuera de él. En
práctica se verifican polos cerrados y diagramas de Nyquist/Bode del modelo
discreto, incluyendo ZOH, retardo y filtros. Analizar solo el controlador
continuo antes de discretizar puede ocultar pérdida de margen de fase.

\subsection{Control deadbeat}
Un diseño deadbeat coloca todos los polos cerrados en $z=0$, de modo que la
respuesta ideal alcanza el valor final en un número finito de muestras. Para
una planta mínima fase puede buscarse $C(z)$ tal que la función cerrada tenga
un denominador monomial. Su atractivo teórico viene acompañado de comandos
grandes, sensibilidad al modelo y mal comportamiento con ceros no mínimos,
retardo o saturación. Por ello no se propone para los MG996R sin identificación
y límites explícitos; se incluye como resultado clásico del control discreto.

\section{Robustez y sistemas no mínimos}
\subsection{Ceros introducidos por muestreo}
La discretización con ZOH puede introducir ceros que no corresponden a ceros
continuos. En plantas de grado relativo alto algunos pueden quedar cerca o
fuera del círculo unitario. Un cero $|z_0|>1$ es no mínimo fase: invertirlo de
forma causal y estable no es posible. Esto limita rapidez de seguimiento y
descarta ciertas estrategias deadbeat o de cancelación.

\subsection{Incertidumbre y estabilidad robusta}
Sea una familia multiplicativa
\begin{equation}
 G_p(z)=G_0(z)[1+W_m(z)\Delta(z)],\qquad
 \|\Delta\|_\infty\le1.
\end{equation}
Una condición de pequeña ganancia suficiente para estabilidad robusta es
\begin{equation}
 \|W_m T\|_\infty<1.
\end{equation}
En Nova, la incertidumbre incluye masa, voltaje, fricción, holgura, ganancia del
servo, retardo y estado de contacto. Antes de aplicar teoría $H_\infty$ formal,
se evaluará una rejilla de modelos y Monte Carlo, conservando las peores
respuestas y no solo el promedio.

\subsection{Perturbaciones y rechazo}
Para perturbación $d$ a la entrada de planta y ruido $n$ en el sensor,
\begin{equation}
 y=T r+G S d-T n.
\end{equation}
Esta igualdad explica el compromiso: ganancia alta reduce $S$ y perturbaciones
de baja frecuencia, pero hace $T\approx1$ y transmite ruido. Contactos e
impactos contienen alta frecuencia; la derivada, el observador y el filtro de
IMU deben diseñarse conjuntamente.

\section{Realización computacional del controlador}
Una función racional debe convertirse en ecuación causal. Para
\begin{equation}
 C(z)=\frac{b_0+b_1z^{-1}+\cdots+b_mz^{-m}}
 {1+a_1z^{-1}+\cdots+a_nz^{-n}},
\end{equation}
la forma directa es
\begin{equation}
 u[k]=-\sum_{i=1}^{n}a_i u[k-i]+
 \sum_{j=0}^{m}b_j e[k-j].
\end{equation}
Cada llamada debe seguir el orden: leer y validar sensores; calcular error;
actualizar estimador; calcular control no saturado; aplicar saturación y
anti-windup; publicar; actualizar memorias; registrar diagnóstico. Las memorias
se inicializan de manera coherente al cambiar de modo para obtener transferencia
sin salto (\emph{bumpless transfer}).

\subsection{Forma directa frente a espacio de estados}
Las formas directas I y II son compactas, pero coeficientes de alto orden pueden
ser sensibles. Se prefieren secciones de segundo orden para filtros o espacio
de estados balanceado para sistemas MIMO. Tras cuantizar coeficientes se vuelven
a calcular polos; estabilidad matemática con doble precisión no garantiza
estabilidad en un microcontrolador de precisión limitada.

\subsection{Pseudocódigo temporal verificable}
\begin{verbatim}
al vencer instante k:
  medir y[k] y su antiguedad
  si dato invalido o supervisor enclavado: entrar en estado seguro
  actualizar x_hat[k]
  obtener referencia nominal q_ref[k]
  limitar correccion residual y tasa de cambio
  calcular u0[k]
  saturar u[k] y actualizar anti-windup
  publicar antes del deadline
  registrar tiempos, estados, saturaciones y modo
\end{verbatim}
Se medirá tiempo de cómputo peor caso, no únicamente promedio. La suma de
adquisición, cómputo y publicación debe dejar margen respecto a $T_s$.

\section{Control predictivo y restricciones}
El control predictivo basado en modelo (MPC) formula, sobre horizonte $N_p$,
\begin{align}
 \min_{\vect u_{0:N_p-1}}\quad &
 \sum_{j=0}^{N_p-1}(\|\vect x_j-\vect x_j^r\|_{\mat Q}^2+
 \|\vect u_j\|_{\mat R}^2)+\|\vect x_{N_p}-\vect x_{N_p}^r\|_{\mat P}^2,\\
 \text{sujeto a}\quad &\vect x_{j+1}=\mat A_d\vect x_j+\mat B_d\vect u_j,\\
 &\vect x_j\in\mathcal X,\quad\vect u_j\in\mathcal U,
 \quad\Delta\vect u_j\in\Delta\mathcal U.
\end{align}
Su ventaja es incorporar límites articulares y de tasa explícitamente. Su costo
es resolver una optimización dentro del deadline y disponer de un modelo
adecuado por contacto. No está implementado ni es requisito actual; se presenta
para completar el panorama de control discreto con restricciones.

\section{Efectos digitales no ideales}
\subsection{Retardo, jitter y pérdida de muestras}
Un retardo entero de $d$ muestras introduce $z^{-d}$. Puede incorporarse como
estados adicionales:
\begin{equation}
 u_1[k+1]=u[k],\quad u_2[k+1]=u_1[k],\ldots,
 \quad u_{aplicado}[k]=u_d[k].
\end{equation}
El retardo agrega fase y reduce margen de estabilidad. El jitter convierte el
sistema en variante en el tiempo; se registrarán periodo medio, desviación,
percentiles y peor caso. Una referencia perdida no debe sustituirse sin
criterio: se define si mantener el último comando, interpolar o entrar en modo
seguro según antigüedad y estado de contacto.

\subsection{Cuantización y precisión numérica}
Para resolución $\Delta$, el cuantizador ideal es
\begin{equation}
 Q(x)=\Delta\operatorname{round}(x/\Delta),\qquad
 e_q=Q(x)-x\in[-\Delta/2,\Delta/2].
\end{equation}
El error no siempre es ruido blanco: puede producir ciclos límite en presencia
de integral y fricción. Se deben documentar resolución PWM, resolución de IMU,
tipos numéricos y conversiones radianes--pulso. La implementación debe evitar
restas catastróficas, acumuladores sin límite y coeficientes mal condicionados.

\subsection{Sistemas multirrate y sincronización}
La marcha, control articular, IMU, contacto y registro pueden tener frecuencias
distintas. Para una señal lenta que alimenta un lazo rápido se usa retención o
interpolación; para diezmar una señal rápida se filtra antes. Cada mensaje debe
llevar sello temporal y marco. La comparación de contactos usa el estado más
reciente, pero su antigüedad forma parte del diagnóstico para no confundir
ausencia física de contacto con pérdida de comunicación.

\section{Control híbrido y discreto de la marcha Nova}
El sistema completo no es solamente LTI. Combina dinámica continua, control
muestreado y eventos de contacto. Se representa como
\begin{equation}
 \mathcal H=(\mathcal X,\mathcal U,\mathcal M,F,G,\mathcal D,\mathcal R),
\end{equation}
donde $\mathcal M$ contiene modos como stand, crawl, step y stop; $F$ describe
la evolución dentro de un modo; $\mathcal D$ las guardas de transición y
$\mathcal R$ los reinicios. El índice de fase obedece
\begin{equation}
 j[k+1]=(j[k]+1)\bmod N,qquad
 n_c[k+1]=n_c[k]+\mathbf 1_{j[k]=N-1}.
\end{equation}
La salida de referencia es una tabla
$\vect q^d[k]=\mathcal Q_m(j[k])$. Para crawl y step, el plan previsto
$\mathcal C_m(j)$ declara una pata en oscilación y tres apoyos.

La arquitectura en cascada propuesta queda
\begin{equation}
 \text{modo/fase}\rightarrow\vect p^d[k]\rightarrow\IK
 \rightarrow\vect q^d[k]\rightarrow
 \text{control articular discreto}\rightarrow\vect\tau[k]
 \rightarrow\text{robot/contacto}.
 \label{eq:cascade}
\end{equation}
El supervisor es una capa paralela de mayor autoridad:
\begin{equation}
 \vect q^{cmd}[k]=
 \begin{cases}
 \vect q^{safe},&s[k]=1,\\
 \vect q^d[k]+\Delta\vect q^{RL}[k],&s[k]=0,
 \end{cases}
\end{equation}
con corrección aprendida limitada por posición y tasa. En el estado actual,
solo altura e inclinación pueden enclavar stand; contacto permanece informativo.

\subsection{Coherencia entre fase prevista y movimiento discreto}
El plan lógico no garantiza el evento físico. Para pata $i$ se definen tiempos
previstos $t_{LO,i}^d,t_{TD,i}^d$ y observados
$t_{LO,i}^o,t_{TD,i}^o$:
\begin{equation}
 \Delta t_{LO,i}=t_{LO,i}^o-t_{LO,i}^d,\qquad
 \Delta t_{TD,i}=t_{TD,i}^o-t_{TD,i}^d.
\end{equation}
La coincidencia ponderada por tiempo es
\begin{equation}
 \gamma_i=\frac{1}{T_E}\int_{T_E}
 \mathbf 1\{c_i^d(t)=c_i^o(t)\}\,dt,
\end{equation}
y la coincidencia simultánea exige igualdad de las cuatro patas. El ensayo
vigente obtuvo $32.550\%$ simultáneo: FL y FR despegaron cerca de
\SI{0.38}{s} tarde y aterrizaron cerca de \SI{1.364}{s} tarde; RL y RR no
despegaron. Por ello, la próxima versión debe modificar transferencia y forma
de oscilación, no simplemente cambiar la etiqueta de contacto.

\begin{workedexample}{cálculo experimental de retardos de contacto}
El ensayo contiene 76 pares válidos de despegue para FL. Para cada ciclo se
calculó
\[
 \Delta t_{LO,FL}^{(r)}=t_{LO,FL}^{o,(r)}-t_{LO,FL}^{d,(r)},
 \qquad r=1,\ldots,76.
\]
La media no se obtuvo restando promedios redondeados, sino promediando los 76
retardos individuales:
\[
 \overline{\Delta t}_{LO,FL}
 =\frac{1}{76}\sum_{r=1}^{76}\Delta t_{LO,FL}^{(r)}
 =\boxed{\SI{0.380296989}{s}}.
\]
Los extremos fueron \SI{0.368741637}{s} y \SI{0.392315495}{s}. Para aterrizaje:
\[
 \overline{\Delta t}_{TD,FL}=\SI{1.364317218}{s}.
\]
Como cada referencia dura \SI{0.18}{s}, estos valores equivalen a
\begin{align*}
 0.380297/0.18&=2.1139\text{ referencias de retraso en despegue},\\
 1.364317/0.18&=7.5795\text{ referencias de retraso en aterrizaje}.
\end{align*}
Para RL y RR no hubo pares de transición: asignar retardo cero habría sido
incorrecto. Se reporta ``sin pares'' porque permanecieron en contacto. Esta
decisión evita convertir ausencia de evento en una medición favorable falsa.
\end{workedexample}

\begin{workedexample}{reproducibilidad de la cadencia corregida}
Los promedios de duración en ciclos 2--13 fueron
$\bar T_1=\SI{4.320013}{s}$ y $\bar T_2=\SI{4.319959}{s}$. Usando el primer
ensayo como referencia:
\begin{align*}
 \Delta_{rel,T}&=100\frac{\bar T_2-\bar T_1}{\bar T_1}\\
 &=100\frac{4.319959-4.320013}{4.320013}
 =\boxed{-0.00125\%}\approx-0.001\%.
\end{align*}
Para avance:
\begin{align*}
 \Delta_{rel,x}&=100\frac{0.023384-0.023338}{0.023338}\\
 &\approx\boxed{0.197\%},
\end{align*}
con pequeñas diferencias frente al $0.194\%$ del informe porque la tabla
impresa redondea las medias antes de mostrarlas; el analizador calcula con
precisión completa. Esta observación explica por qué los resultados finales
deben obtenerse desde CSV, no desde números copiados de una tabla redondeada.
\end{workedexample}

\section{Procedimiento completo de diseño discreto}
Para que el capítulo se traduzca en ingeniería reproducible se seguirá esta
secuencia:
\begin{enumerate}
 \item definir variables, entradas, salidas, marcos, restricciones y punto de
 operación;
 \item fijar especificaciones: error, sobreimpulso, establecimiento, esfuerzo,
 robustez, contacto y seguridad;
 \item identificar o derivar el modelo continuo y cuantificar incertidumbre;
 \item elegir $T_s$ con base en ancho de banda, sensores, cómputo y retardo;
 \item discretizar con ZOH exacto y comparar contra Tustin o aproximaciones;
 \item verificar polos, ceros, estabilidad, controlabilidad y observabilidad;
 \item seleccionar PID, realimentación de estados, LQR u otro método coherente;
 \item añadir observador, integral, saturaciones, anti-windup y retardos;
 \item simular casos nominales y extremos, incluyendo pérdida de datos;
 \item implementar con sellos temporales, límites y estado seguro;
 \item verificar ecuaciones contra una implementación independiente;
 \item validar primero en MuJoCo y Gazebo y después progresivamente en hardware.
\end{enumerate}

\subsection{Especificaciones mínimas antes de sintonizar}
No se aceptará una ganancia solo porque ``se ve estable''. Para cada lazo se
registrarán: $T_s$, ancho de banda o polos deseados, sobreimpulso, tiempo de
establecimiento, error RMS y máximo, saturación, energía, ruido, retardo máximo,
margen de ganancia/fase cuando aplique y región de posturas/contactos válida.
Se separarán datos de identificación, sintonización y validación.

\subsection{Ejemplo numérico didáctico de un servo de primer orden}
Como ejemplo, no como identificación del MG996R, sea
\begin{equation}
 G(s)=\frac{1}{\tau s+1},\qquad \tau=\SI{0.12}{s},\quad T_s=\SI{0.02}{s}.
\end{equation}
La discretización ZOH exacta produce
\begin{equation}
 x[k+1]=e^{-T_s/\tau}x[k]+(1-e^{-T_s/\tau})u[k]
 =0.8465x[k]+0.1535u[k].
\end{equation}
Su polo $z=0.8465$ está dentro del círculo unitario. Con control proporcional
$u=K_p(r-x)$, el polo cerrado es
\begin{equation}
 z_{cl}=0.8465-0.1535K_p.
\end{equation}
Exigir $|z_{cl}|<1$ da aproximadamente $-1.00<K_p<12.03$. Este intervalo solo
garantiza estabilidad del modelo ideal; restricciones, ruido y retardo reducen
el rango útil. El ejemplo muestra el procedimiento: modelo, discretización,
polo, desigualdad y luego validación.

\section{Qué está implementado y qué es diseño propuesto}
\begin{longtable}{p{5.0cm}p{4.0cm}p{5.0cm}}
\toprule Elemento&Estado&Observación\\\midrule
Secuenciador de marcha discreto&Implementado&24 o 32 referencias; vencimiento
sin deriva acumulativa.\\
Publicación de fase y contacto previsto&Implementado&JSON sincronizado con cada
referencia.\\
Interpolación articular&Implementada en \texttt{ros2\_control}&Debe caracterizarse
su equivalente discreto exacto.\\
Supervisor de altura e inclinación&Implementado, provisional&Ordena stand y se
enclava; no sustituye parada física.\\
PID propio de los MG996R&No implementado&El servo tiene electrónica interna no
identificada; la interfaz PCA9685 sigue pendiente.\\
Realimentación de estados, LQR y observador&Formulación propuesta&Requieren
modelo identificado y sensores suficientes.\\
IMU y contacto equivalente en MuJoCo&Pendiente&Necesarios para observación y
comparación homogénea.\\
Polígono y margen en línea&Pendiente&Debe usar contactos físicamente válidos.\\
Corrección RL acotada&Pendiente&Será residual, nunca PWM directo.\\
\bottomrule
\end{longtable}

\section{Resultados nominales}
Para $\vect q_i=(0.10,0.42,-0.84)\trans$ rad:
\begin{align}
\mat M_i&=\begin{bmatrix}
0.00904783&0.00035891&-0.00041176\\
0.00035891&0.00673928&0.00195619\\
-0.00041176&0.00195619&0.00316985
\end{bmatrix}\ \si{kg.m^2},\\
\operatorname{eig}(\mat M_i)&=(0.00226642,0.00758520,0.00910534)\ \si{kg.m^2},\\
\vect g_i&=(0.30982,0.03891,-0.04464)\trans\ \si{N.m}.
\end{align}
\begin{table}[ht]\centering\caption{Indicadores nominales.}
\begin{tabular}{ll}\toprule Indicador&Resultado\\\midrule
COM corporal&$(-0.00586,0,-0.03882)$ m\\
Margen estático con cuatro apoyos&\SI{0.07313}{m}\\
Mayor par gravitacional&\SI{0.30982}{N.m}\\
Proporción del par de bloqueo&$28.7\%$\\\bottomrule
\end{tabular}\end{table}

\section{Implementación y verificación}
La implementación se encuentra en
\texttt{nova\_gait\_controller/mathematical\_model.py}. Se comprobaron FK--IK,
Jacobiano frente a diferencias centrales, potencia virtual, simetría y
definitud positiva de $\mat M$, Coriolis nulo a velocidad cero, dinámica
inversa finita, COM, contacto, margen estático, actuador y consistencia entre
código, URDF/Xacro y MuJoCo. La verificación actual registra 33 pruebas aprobadas,
URDF válido y paquetes ROS 2 compilados.

El ejemplo reproducible se ejecuta con:
\begin{verbatim}
cd /home/pavilion/Documentos/Cuadrupedo
source /opt/ros/humble/setup.bash
source install/setup.bash
python3 scripts/ejemplo_modelo_matematico.py
\end{verbatim}
Para $\vect q=(0.10,0.42,-0.84)$ rad,
$\dot{\vect q}=(0.10,-0.05,0.08)$ rad/s,
$\ddot{\vect q}=(0.20,0.10,-0.15)$ rad/s$^2$ y
$\vect f=(1,0.5,5)$ N, reconstruye término por término \cref{eq:inverse-dynamics}.

\section{Identificación paramétrica}
Los parámetros se separan en geometría, inercia, actuadores y contacto. Cada
registro debe contener símbolo, unidad, valor nominal, método, instrumento,
repeticiones, estimación, incertidumbre, fecha, condiciones y fuente. El ajuste
ponderado se formula como
\begin{equation}
 \vect\theta^*=\argmin_{\vect\theta}\sum_{k=1}^N
 \|\vect y_k-\hat{\vect y}_k(\vect\theta)\|_{\mat W}^2.\label{eq:identification}
\end{equation}
Los datos de ajuste deben separarse de los de validación. El orden seguro es:
medir sin energizar; caracterizar un MG996R asegurado con límite de corriente;
dimensionar potencia y parada; calibrar unidades; y solo entonces probar patas.

\section{Validación cuantitativa}
Se distinguen verificación interna, validación en simulación y validación
física. Para posición de pie:
\begin{align}
 \vect e_p(k)&=\vect p_{\mathrm{ref}}(k)-\vect p_{\mathrm{modelo}}(k),\\
 \operatorname{RMSE}_p&=\sqrt{\frac1N\sum_{k=1}^N\|\vect e_p(k)\|^2}.
\end{align}
También se reportan error máximo, retardo, altura, roll, pitch, avance, deriva y,
si existe instrumentación, fuerza, par y corriente. La línea base de gateo ya
fue registrada y repetida; la cadencia corregida mostró diferencias menores al
0.2\% en las medias. La marcha paso fue validada en Gazebo y articularmente en
MuJoCo. Un ensayo adicional de 76 ciclos cuantificó solo 32.550\% de
coincidencia entre contacto previsto y observado y ausencia de despegue trasero.
La trayectoria debe corregirse antes de interpretar el margen de soporte. El
hardware requiere medición física real.

\section{Sensibilidad e incertidumbre}
\begin{align}
 S_{y,\theta_j}&=\frac{\partial y}{\partial\theta_j},&
 \widetilde S_{y,\theta_j}&=\frac{\theta_j}{y}
 \frac{\partial y}{\partial\theta_j},\\
 \frac{\partial y}{\partial\theta_j}&\approx
 \frac{y(\theta_j+h)-y(\theta_j-h)}{2h}.
\end{align}
Se variarán longitudes, masas, COM, fricción, tensión, ganancias, retardos y
contacto. Las salidas son posición de pie, margen, par, altura, inclinación y
avance. Monte Carlo complementará la sensibilidad local. Los parámetros con
mayor efecto reciben prioridad experimental.

\section{Consistencia entre modelo, código y simuladores}
Una fuente de parámetros debe alimentar Python, URDF/Xacro y MJCF. Las pruebas
comparan geometría, masas, tensores, límites, nombres, ejes, FK/IK, Jacobiano,
COM y saturaciones. Se documentan versiones, solver, paso de integración,
unidades y marcos. La ganancia observada en Gazebo Humble es $0.1$ aunque el
YAML contenga $0.5$; por ello $0.5$ no es un valor validado.

\section{Evidencia requerida y criterio de cierre}
La versión final deberá incorporar: robot con marcos y dimensiones; pata con
ángulos; workspace 3D; mapa de $\sigma_{\min}$ o $\kappa$; trayectoria $x$--$z$
con velocidad y aceleración; fases temporales; polígono de soporte y COM; y
curvas sincronizadas modelo--Gazebo--MuJoCo--experimento.

El modelo estará cerrado cuando pueda seguirse sin saltos la cadena desde
trayectoria hasta coordenadas, velocidades, aceleraciones, pares, contacto y
respuesta corporal; cuando alcance, singularidad y límites estén definidos; y
cuando las predicciones se comparen con datos independientes. Las gráficas
reproducibles de workspace, singularidad y velocidades articulares ya se
incorporaron al documento final de tesis. Permanecen pendientes la
implementación verificada de $\dot{\mat J}$, la identificación física y la
validación experimental.

\section{Registro de procedencia y uso de las ecuaciones}
Esta tabla cierra la trazabilidad a nivel de ecuación. ``Procedencia'' significa
la teoría de la que se deriva la forma matemática; ``adaptación Nova'' identifica
qué se decidió específicamente para este robot. Las referencias T01--T11 están
definidas en la matriz inicial. Cuando varias ecuaciones aparecen juntas es
porque constituyen una sola derivación y no deben aplicarse aisladamente.

\scriptsize
\begin{longtable}{p{2.5cm}p{2.1cm}p{3.1cm}p{3.5cm}p{3.3cm}}
\caption{Trazabilidad ecuación--método--propósito.}\label{tab:eqtrace}\\
\toprule Ecuación o familia&Procedencia&Variables que entrega&Adaptación Nova&Uso verificable en la tesis\\
\midrule\endfirsthead
\toprule Ecuación o familia&Procedencia&Variables que entrega&Adaptación Nova&Uso verificable en la tesis\\
\midrule\endhead
\cref{eq:qlevels}&Definición propia&$\vect q_i$&Orden coxa--fémur--tibia y
orden FL--FR--RL--RR&Evitar permutaciones entre URDF, ROS y vectores del modelo.\\
\cref{eq:homogeneous,eq:chain,eq:worldfoot}&T01,T02&$\mat T$, $\mat R$,
$\vect p_i$&Ejes REP--103, caderas $d_x,d_y$ y cadena Nova&Transformar pie
local a cuerpo y mundo; contraste FK--URDF.\\
\cref{eq:fkx,eq:fkzp,eq:fky,eq:fkz}&T02,T03&$x,y,z$&Longitudes Nova,
$s_i$ y ceros articulares propios&Generar referencias cartesianas y comprobar
posición del pie.\\
\cref{eq:ikzp,eq:ikcoxa,eq:D,eq:kneebranch,eq:ikfemur,eq:iklimits}&T02,T03&
$q_c,q_f,q_t,D$&Rama de rodilla negativa, límites URDF y continuidad&Convertir
cada punto de marcha en una orden articular segura.\\
\cref{eq:workspace}&T02&$\mathcal W_i$&Caja articular y geometría Nova&Rechazar
objetivos inalcanzables antes de publicarlos.\\
\cref{eq:jacdef,eq:jacobian,eq:singularity}&T02&$\mat J_i,w_i$&Derivación
analítica de la FK Nova&Velocidad, fuerza--par y proximidad a singularidad.\\
\cref{eq:accel,eq:jdot}&T02&$\ddot{\vect p},\dot{\mat J}$&Diferencia central
con paso declarado&Obtener aceleración cartesiana y verificar derivadas.\\
\cref{eq:dynamics,eq:kinetic,eq:potential,eq:mg,eq:christoffel}&T04&
$\mat M,\vect c,\vect g$&Masas, COM e inercias nominales Nova&Separar los
términos físicos del par y preparar identificación.\\
\cref{eq:inverse-dynamics}&T04&$\vect\tau$&Fricción y fuerzas externas con
signos Nova&Comparar demanda calculada con capacidad del servo.\\
\cref{eq:actuator}&T06&$\tau_{\max}$&Escalado nominal con tensión y velocidad&
Filtro preliminar de factibilidad; no reemplaza ensayo de banco.\\
\cref{eq:complementarity,eq:frictioncone}&T05&$\lambda_n,\vect\lambda_t$&
$\mu=0.90$ nominal y normal vertical&Decidir apoyo y detectar deslizamiento en
simulación.\\
\cref{eq:staticmargin}&T05&$\vect r_{COM},m_s$&Proyección sobre casco convexo
de pies activos&Cuantificar estabilidad estática con signo y unidad.\\
\cref{eq:gaitphase,eq:stance,eq:swing,eq:motionchain}&T02,T03 y adaptación
propia&$\phi_i,x,z,\vect q_i$&Orden FL--RR--FR--RL, $L,H,N,T$ Nova&Construir
la marcha y conservar la cadena fase--contacto.\\
\cref{eq:quintic}&T07&$\vect p(t)$ y derivadas&Interpolación entre puntos Nova&
Evitar saltos de velocidad y aceleración en trayectorias futuras.\\
\cref{eq:deadline,eq:sampling}&T07,T10&$k,t_k,T_s$&Reloj monotónico y avance
por vencimiento&Impedir deriva acumulativa y hacer repetible la cadencia.\\
\cref{eq:nyquist,eq:zoh,eq:difference,eq:ztransform,eq:pulse-transfer,eq:szmap}&
T07&$x[k],X(z),G(z)$&Aplicación al periodo real del nodo&Pasar correctamente
del fenómeno continuo al algoritmo digital.\\
\cref{eq:sscontinuous,eq:ssdiscrete,eq:exactdisc,eq:blockexp}&T04,T07&
$\mat A_d,\mat B_d$&ZOH exacto calculado por exponencial en bloques&Predecir
estado por muestra y habilitar análisis discreto.\\
\cref{eq:linearjoint,eq:jointss}&T04,T07&$J,b,k_t,\vect x$&Eje de fémur en
postura nominal&Ejemplo reproducible que une dinámica articular y control.\\
\cref{eq:discretesolution,eq:schur,eq:jury2,eq:dlyap,eq:finalvalue}&T07&
polos, $\mat P$, error final&Aplicación al modelo discretizado Nova&Probar
estabilidad y desempeño, no juzgarlos solo por una gráfica.\\
\cref{eq:ctrb,eq:obsv}&T07&$\mathcal C,\mathcal O$&Estados y sensores que
realmente estarán disponibles&Comprobar si se puede controlar y reconstruir el
estado propuesto.\\
\cref{eq:pid,eq:antiwindup}&T07&$u[k],I[k]$&Límites de referencia/par Nova&
Plantear un PID implementable sin acumulación integral en saturación.\\
\cref{eq:statefeedback,eq:augmented,eq:lqrcost}&T08&$\mat K,\vect x_I$&Pesos
y modelo nominal declarados&Diseñar y comparar realimentación con seguimiento
sin ocultar la sintonización.\\
\cref{eq:observer}&T07,T09&$\hat{\vect x}$&Salidas disponibles aún por fijar&
Propuesta para estimar velocidad/estado; pendiente de sensores identificados.\\
\cref{eq:cascade}&T07 y arquitectura propia&$\vect p_d,\vect q_d,\vect\tau_d$&
Capas cartesiana, articular y servo Nova&Delimitar responsabilidades y ubicar
qué señales deben medirse en cada lazo.\\
\cref{eq:identification}&Método de mínimos cuadrados ponderados&$\vect\theta^*$&
Pesos, repeticiones y partición de datos por definir&Reemplazar parámetros
nominales por estimaciones físicas con incertidumbre.\\
\bottomrule
\end{longtable}
\normalsize

\section{Fuentes técnicas}
Las fuentes web siguientes se verificaron el 16 de agosto de 2026. Se citan
como base de método o contrato de software; los resultados numéricos Nova
proceden del código y los experimentos identificados en la matriz de
trazabilidad.
\begin{enumerate}
 \item ROS Enhancement Proposal 103, convenciones de coordenadas:
 \url{https://www.ros.org/reps/rep-0103.html}.
 \item NovaSM3, ``The Muscles'': \url{https://novaspotmicro.com/the-muscles.html}.
 \item Repositorio Nova-SM3: \url{https://github.com/cguweb-com/Arduino-Projects/tree/main/Nova-SM3}.
 \item Lista NovaSM3 v5.2b: \url{https://novaspotmicro.com/parts-list.html}.
 \item TowerPro, MG996R: \url{https://towerpro.com.tw/product/mg996R/}.
 \item DSServo, DS3218: \url{https://www.dsservo.com/down.asp?id=22}.
 \item Mike4192, Spot Micro Kinematics: \url{https://github.com/mike4192/spot_micro_kinematics_python}.
 \item K. Ogata, \emph{Discrete-Time Control Systems}, 2.a ed., Prentice Hall, 1995.
 \item G. F. Franklin, J. D. Powell y M. L. Workman, \emph{Digital Control of
 Dynamic Systems}, 3.a ed., Addison-Wesley, 1998.
 \item K. J. \AA str\"om y B. Wittenmark, \emph{Computer-Controlled Systems},
 3.a ed., Prentice Hall, 1997.
 \item C.-T. Chen, \emph{Linear System Theory and Design}, 3.a ed., Oxford
 University Press, 1999.
 \item R. Tedrake, MIT Underactuated Robotics, ``Multi-Body Dynamics'':
 \url{https://underactuated.mit.edu/multibody.html}.
 \item R. Tedrake, MIT Underactuated Robotics, ``Linear Quadratic Regulators'':
 \url{https://underactuated.mit.edu/lqr.html}.
 \item R. E. Kalman, ``A New Approach to Linear Filtering and Prediction
 Problems'', \emph{Journal of Basic Engineering}, 82(1), 35--45, 1960,
 \url{https://doi.org/10.1115/1.3662552}.
 \item ROS 2, definición oficial de \texttt{JointTrajectory}:
 \url{https://docs.ros.org/en/jazzy/p/trajectory_msgs/msg/JointTrajectory.html}.
 \item ROS 2 Control, documentación de \texttt{joint\_trajectory\_controller}:
 \url{https://docs.ros.org/en/ros2_packages/rolling/api/joint_trajectory_controller/doc/userdoc.html}.
 \item Gazebo Sim, tutorial oficial de sensores IMU y contacto:
 \url{https://gazebosim.org/docs/garden/sensors/}.
\end{enumerate}

\section{Conclusión}
Se consolidó una estructura nominal computable que enlaza geometría,
cinemática, dinámica, actuadores, contacto, estabilidad, generación de marcha y
control discreto. Se desarrollaron muestreo, transformada Z, discretización,
estabilidad, controlabilidad, observabilidad, PID, realimentación de estados,
LQR, observadores, robustez, implementación digital y sistemas híbridos. Las
formulaciones de control avanzado se presentan como diseño verificable, no como
funciones ya implementadas en el robot.
Esta versión reconstruye todas las expresiones como matemáticas LaTeX nativas y
define workspace, singularidades, aceleraciones, complementariedad,
identificación, validación, sensibilidad y trazabilidad. Solo después de
sustituir parámetros estimados y cuantificar errores frente al robot podrá
considerarse representativa del ejemplar físico.

\end{document}
